% Options for packages loaded elsewhere
\PassOptionsToPackage{unicode}{hyperref}
\PassOptionsToPackage{hyphens}{url}
\documentclass[
  12pt,
]{article}
\usepackage{xcolor}
\usepackage[margin=1in]{geometry}
\usepackage{amsmath,amssymb}
\setcounter{secnumdepth}{-\maxdimen} % remove section numbering
\usepackage{iftex}
\ifPDFTeX
  \usepackage[T1]{fontenc}
  \usepackage[utf8]{inputenc}
  \usepackage{textcomp} % provide euro and other symbols
\else % if luatex or xetex
  \usepackage{unicode-math} % this also loads fontspec
  \defaultfontfeatures{Scale=MatchLowercase}
  \defaultfontfeatures[\rmfamily]{Ligatures=TeX,Scale=1}
\fi
\usepackage{lmodern}
\ifPDFTeX\else
  % xetex/luatex font selection
  \setmainfont[]{Times New Roman}
  \setmonofont[]{Consolas}
\fi
% Use upquote if available, for straight quotes in verbatim environments
\IfFileExists{upquote.sty}{\usepackage{upquote}}{}
\IfFileExists{microtype.sty}{% use microtype if available
  \usepackage[]{microtype}
  \UseMicrotypeSet[protrusion]{basicmath} % disable protrusion for tt fonts
}{}
\usepackage{setspace}
\makeatletter
\@ifundefined{KOMAClassName}{% if non-KOMA class
  \IfFileExists{parskip.sty}{%
    \usepackage{parskip}
  }{% else
    \setlength{\parindent}{0pt}
    \setlength{\parskip}{6pt plus 2pt minus 1pt}}
}{% if KOMA class
  \KOMAoptions{parskip=half}}
\makeatother
\usepackage{longtable,booktabs,array}
\usepackage{caption}
\captionsetup[table]{skip=6pt}
\newcounter{none} % for unnumbered tables
\usepackage{calc} % for calculating minipage widths
% Correct order of tables after \paragraph or \subparagraph
\usepackage{etoolbox}
\makeatletter
\patchcmd\longtable{\par}{\if@noskipsec\mbox{}\fi\par}{}{}
\makeatother
\setlength{\emergencystretch}{3em} % prevent overfull lines
\providecommand{\tightlist}{%
  \setlength{\itemsep}{0pt}\setlength{\parskip}{0pt}}
\usepackage{newunicodechar}
\newunicodechar{₀}{\textsubscript{0}}
\newunicodechar{₁}{\textsubscript{1}}
\newunicodechar{₂}{\textsubscript{2}}
\newunicodechar{₃}{\textsubscript{3}}
\newunicodechar{ᵢ}{\textsubscript{i}}
\newunicodechar{ⱼ}{\textsubscript{j}}
\newunicodechar{ₖ}{\textsubscript{k}}
\newunicodechar{σ}{\ensuremath{\sigma}}
\newunicodechar{μ}{\ensuremath{\mu}}
\newunicodechar{ε}{\ensuremath{\varepsilon}}
\newunicodechar{α}{\ensuremath{\alpha}}
\newunicodechar{β}{\ensuremath{\beta}}
\newunicodechar{γ}{\ensuremath{\gamma}}
\newunicodechar{λ}{\ensuremath{\lambda}}
% APA 7: texto alineado a la izquierda (flush left, ragged right)
\usepackage{ragged2e}
\raggedright
\setlength{\emergencystretch}{3em}
% APA 7: interlineado doble
\usepackage{setspace}
\doublespacing
% APA 7: sangría de 0.5 pulgadas en primera línea de párrafo
\setlength{\parindent}{0.5in}
\setlength{\parskip}{0pt}
\usepackage{bookmark}
\IfFileExists{xurl.sty}{\usepackage{xurl}}{} % add URL line breaks if available
\urlstyle{same}
% fallback for those not using the hyperref driver hyperxmp:
\makeatletter
\@ifundefined{xmpquote}{\newcommand{\xmpquote}[1]{#1}}{}
\makeatother
\hypersetup{
  hidelinks,
  pdfcreator={LaTeX via pandoc}}

\author{}
\date{}

\begin{document}

\setstretch{2}
\section{1. INTRODUCCIÓN}\label{introducciuxf3n}

\subsection{1.1. Realidad Problemática}\label{realidad-problemuxe1tica}

La proliferación masiva de dispositivos de Internet de las Cosas (IoT)
en entornos interiores —edificios inteligentes, plantas industriales,
campus universitarios y instalaciones sanitarias— ha generado un
escenario de alta densidad de nodos inalámbricos que compiten por
recursos de comunicaciones finitos y, en muchos casos, degradados por
las condiciones del canal. La transmisión de señales de radiofrecuencia
en interiores se caracteriza por fenómenos físicos complejos como la
multipropagación, el desvanecimiento selectivo en frecuencia, la
atenuación por materiales de construcción y la interferencia co-canal,
los cuales distorsionan de manera significativa el enlace de
comunicaciones y elevan la tasa de error a nivel físico. Este
comportamiento del medio de propagación impone un límite teórico y
práctico sobre la calidad de servicio (QoS) que puede ofrecerse a las
aplicaciones críticas IoT, tales como monitoreo de procesos, control de
seguridad y telemedicina, que requieren latencia acotada, confiabilidad
elevada y disponibilidad continua.

En paralelo, la congestión dinámica del medio se intensifica a medida
que la densidad de nodos aumenta. Los protocolos de acceso al medio de
bajo consumo, ampliamente adoptados en el ecosistema IoT (por ejemplo,
IEEE 802.15.4/Zigbee, LoRaWAN en variantes indoor y redes mesh sobre
802.11), operan en bandas de espectro compartidas y no licenciadas
(ISM), donde la probabilidad de colisión de tramas crece de forma no
lineal con el número de emisores activos por unidad de área. El problema
se agrava por el fenómeno del \emph{hidden node}, las variaciones
temporales del tráfico (patrones estocásticos de generación de paquetes)
y la ausencia de un coordinador central capaz de reasignar recursos de
forma óptima en tiempo real. De este modo, se establece un conflicto
estructural entre la capa física y la capa MAC: mientras la propagación
multipath reduce la robustez del enlace, la congestión dinámica degrada
la eficiencia espectral y eleva la energía disipada por reintentos,
configurando un círculo vicioso de rendimiento que limita el escalado
sostenible de redes IoT densas.

Las soluciones convencionales para mitigar este conflicto se han basado
históricamente en infraestructuras de alto costo: despliegue de sistemas
DAS (\emph{Distributed Antenna Systems}) con fibra óptica hasta cada
punto de acceso, estaciones base de macrocelda reforzadas, gateways
industriales certificados con antenas direccionales, o la adquisición de
espectro licenciado con garantías regulatorias. Si bien estas
arquitecturas demostraron eficacia en contextos de telecomunicaciones
celulares tradicionales, resultan económicamente inviables para la gran
mayoría de las aplicaciones IoT de interiores, cuyo modelo de negocio
exige costos de despliegue por nodo en el rango de unidades a decenas de
dólares, consumo energético compatible con baterías de larga duración y
administración descentralizada sin depender de personal especializado
permanente. La inversión en infraestructura centralizada además
concentra un punto único de falla, contradiciendo los principios de
resiliencia que la propia arquitectura IoT promete.

En consecuencia, existe una brecha tecnológica significativa entre las
soluciones disponibles y las necesidades reales de las redes IoT
interiores: se requieren mecanismos adaptativos, de bajo costo y
distribuidos, capaces de compensar simultáneamente los efectos de la
propagación física y la congestión dinámica del medio. Abordar esta
problemática demanda un enfoque que integre el modelado del canal indoor
con estrategias de gestión inteligente del espectro y del acceso al
medio, orientado a maximizar la eficiencia de la red sin recurrir a
infraestructura capital intensiva. La presente investigación se enmarca
en esta necesidad, proponiendo soluciones viables desde el punto de
vista técnico y económico para la próxima generación de redes IoT densas
en entornos interiores.

{[}INSERTAR AQUÍ TABLA: Comparativa de soluciones al conflicto
físico–MAC en IoT interiores — infraestructura centralizada de alto
costo vs.~propuesta distribuida de bajo costo{]}

\subsection{1.2. Preguntas de
Investigación}\label{preguntas-de-investigaciuxf3n}

La realidad problemática descrita evidencia que la calidad de servicio
(QoS) de una red inalámbrica interior no es una propiedad atribuible a
un solo componente del sistema, sino un resultado \emph{emergente} de la
interacción entre las condiciones físicas del canal (atenuación,
multipropagación, interferencia) y el comportamiento dinámico de la pila
de protocolos (acceso al medio, colisiones, reintentos, sobrecarga de
confirmaciones). Preguntar por el QoS en interiores equivale, por tanto,
a preguntar por una relación funcional entre conjuntos de variables de
naturaleza distinta —físicas y de red— cuya dependencia raramente es
lineal y que, en consecuencia, exige un instrumento de medición capaz de
registrar ambas familias de variables de forma simultánea, sincronizada
y reproducible. Ese instrumento es el prototipo IoT híbrido que se
propone diseñar: nodos de bajo costo basados en ESP32 que transmiten por
Wi-Fi IEEE 802.11 publicaciones MQTT y que, al mismo tiempo, reportan
los indicadores de radiofrecuencia del enlace.

En tal virtud, la formulación de las preguntas de investigación se
estructura como un sistema jerárquico y cerrado: una pregunta general
que encuadra el problema sistémico y cuatro preguntas específicas que,
en correspondencia con los cuatro eslabones de la cadena experimental
(canal de propagación, comportamiento del QoS, correlación
estadístico-aprendizaje y viabilidad del instrumento), permiten
desagregar la respuesta en problemas de ingeniería abordables,
verificables y compatibles con el diseño factorial 3×3×3 y con el umbral
de desempeño del 85\% que se asume como criterio de éxito.

{[}INSERTAR AQUÍ FLUJOGRAMA DE: Sistema jerárquico de preguntas de
investigación y su articulación con la cadena experimental (PE1 canal →
PE2 red → PE3 Machine Learning → PE4 viabilidad){]}

\subsubsection{1.2.1. Pregunta General}\label{pregunta-general}

\textbf{¿En qué medida un prototipo IoT híbrido basado en
microcontroladores ESP32, enlace Wi-Fi IEEE 802.11 y publicaciones MQTT
permite evaluar la calidad de servicio de las redes inalámbricas
interiores mediante la correlación de variables físicas de
radiofrecuencia con variables dinámicas de red a través de técnicas de
Machine Learning, y en qué forma dicho instrumento permite verificar de
manera reproducible el umbral de desempeño del 85\% fijado como criterio
de éxito?}

\emph{Trasfondo de ingeniería.} La pregunta se justifica en la
naturaleza \emph{transversal} (cross-layer) del fenómeno que se pretende
medir. En telecomunicaciones, la optimización o la evaluación de una
sola capa presenta rendimientos decrecientes cuando el cuello de botella
emerge de la interacción entre capas: ninguna métrica de red describe
por sí sola la degradación causada por un canal con desvanecimiento
selectivo, y ninguna medición de potencia de recepción anticipa por sí
sola la latencia generada por las colisiones del mecanismo CSMA/CA. Para
caracterizar el QoS es imprescindible observar, en la misma escala de
tiempo, los indicadores físicos (RSSI, SNR, ruido de fondo) y los
indicadores de red (latencia extremo a extremo, jitter, razón de entrega
de tramas, pérdida de paquetes y buen uso). El prototipo propuesto se
convierte, así, en el \emph{instrumento de medición} que hace observable
esa relación, y la pregunta general encuadra la investigación como un
problema de \emph{ingeniería de sistemas de comunicaciones y de
instrumentación}, donde la variable de diseño no es un componente
aislado, sino el conjunto hardware–software–algoritmo que gobierna la
adquisición, el tratamiento y la interpretación conjunta de ambos tipos
de señales bajo restricciones económicas explícitas de despliegue
masivo.

\subsubsection{1.2.2. Preguntas
Específicas}\label{preguntas-especuxedficas}

\paragraph{Pregunta Específica 1
(PE1)}\label{pregunta-especuxedfica-1-pe1}

\textbf{¿Qué modelo de propagación en interiores, calibrado para la
banda ISM de 2,4 GHz y las geometrías típicas de instalaciones densas,
describe con suficiente fidelidad las variables físicas de
radiofrecuencia —RSSI, SNR, potencia de recepción y ruido de fondo— que
el prototipo ESP32 es capaz de estimar en campo, y qué relación guardan
dichas variables con la atenuación por materiales de construcción y con
los mecanismos de multipropagación presentes en el escenario de
medición?}

\emph{Trasfondo de ingeniería.} Esta pregunta sitúa la investigación en
el dominio de la \emph{propagación de radiofrecuencia y caracterización
de canales}. El canal de comunicación indoor se modela, según la
complejidad requerida, mediante modelos deterministas (\emph{ray
tracing} de alta frecuencia), semi-empíricos (modelos de pérdida de
distancia con exponente de atenuación ajustable, el modelo de la
Recomendación ITU-R P.1238 de la forma L = 20 log₁₀ f + N log₁₀ d +
Lf(n) − 28 dB, o modelos de múltiples paredes tipo Motley-Keenan) o
estocásticos (distribuciones Rayleigh, Rician o Nakagami-\emph{m} para
el desvanecimiento de amplitud, con desvanecimiento por sombra de
distribución log-normal). Cada modelo implica una relación distinta
entre coste computacional y fidelidad, y su elección determina la
validez de todas las evaluaciones posteriores: sin una caracterización
correcta del exponente de pérdida de distancia, del desvanecimiento por
sombra y del \emph{delay spread} propio del escenario, cualquier
conclusión sobre QoS se formularía sobre una abstracción irreal del
medio y los resultados no serían trasladables a un despliegue físico
distinto. Adicionalmente, la pregunta exige identificar los parámetros
\emph{medibles en campo} con el hardware de bajo costo disponible,
porque condiciona toda la viabilidad del prototipo: un ESP32 reporta el
RSSI desde el controlador de interfaz Wi-Fi y permite estimar el ruido
de fondo y la relación señal-ruido, pero no la matriz de canal compleja
completa. En consecuencia, PE1 debe resolver dos problemas de ingeniería
concretos: (i) ajustar por mínimos cuadrados los parámetros del modelo
de pérdida de distancia a partir de la campaña de medición propia,
verificando su bondad de ajuste mediante el error medio cuadrático y el
coeficiente de determinación; y (ii) establecer el rango de operación
del instrumento en el cual el RSSI reportado conserva monotonicidad
respecto de la distancia, dada la conocida no linealidad y el efecto de
las reflexiones constructivas y destructivas de la multipropagación. En
síntesis, PE1 establece el \emph{cimiento empírico-matemático} sobre el
cual se construye la familia de variables explicativas de todo el
análisis posterior.

{[}INSERTAR AQUÍ TABLA: Familias de modelos de propagación en interiores
— deterministas, semi-empíricos (ITU-R P.1238, Motley-Keenan) y
estocásticos (Rayleigh, Rician, Nakagami-m) — y su adecuación al
instrumento ESP32{]}

\paragraph{Pregunta Específica 2
(PE2)}\label{pregunta-especuxedfica-2-pe2}

\textbf{¿Cómo se comportan las métricas de calidad de servicio —latencia
extremo a extremo, jitter, razón de entrega de tramas (PDR), pérdida de
paquetes y buen uso— de un enlace MQTT sobre Wi-Fi IEEE 802.11 en
interiores, al variar de forma controlada el escenario de propagación,
la carga de ofrecimiento y el nivel de QoS del protocolo, y en qué
combinaciones del diseño factorial experimental 3×3×3 se alcanza o no el
umbral del 85\% fijado como criterio de éxito?}

\emph{Trasfondo de ingeniería.} PE2 corresponde al campo de la
\emph{evaluaración experimental de protocolos y de la ingeniería de
calidad de servicio}. La definición de QoS de la Recomendación ITU-T
E.800 —conjunto de efectos de un servicio que condicionan la
satisfacción del usuario— debe operacionalizarse en magnitudes medibles
antes de poder verificarse: la latencia extremo a extremo (con la
referencia de 150 ms de la Recomendación ITU-T G.114 como límite
prácticamente imperceptible para servicios interactivos), la variación
de retardo o jitter, la razón de entrega de tramas, el porcentaje de
pérdida y el buen uso en bits por segundo. Sobre el plano de enlace,
Wi-Fi IEEE 802.11 administra el medio mediante DCF/CSMA/CA con retroceso
exponencial binario y ventana de contienda, mecanismo cuyo
comportamiento se modela con cadenas de Markov —el paradigma clásico del
modelo de Bianchi— y que degrada de forma no lineal con la carga y con
las interferencias co-canal y de canal adyacente propias de la banda ISM
de 2,4 GHz, compartida con Bluetooth, Zigbee y equipos vecinos no
coordinados. Sobre el plano de aplicación, MQTT introduce una dimensión
adicional que suele subestimarse: sus tres niveles de servicio (QoS 0
sin confirmación, QoS 1 con confirmación de recepción y QoS 2 con ciclo
completo de cuatro manos) intercambian confiabilidad por sobrecarga de
tráfico y por latencia, de modo que el nivel elegido modifica
simultáneamente el tiempo de entrega y la probabilidad de duplicación o
pérdida. Finalmente, la pregunta incorpora el \emph{diseño factorial} no
como un recurso estadístico ornamental, sino como la única forma de
estimar simultáneamente los \emph{efectos principales} de los tres
factores y sus \emph{interacciones} —por ejemplo, si el efecto de la
carga de tráfico sobre la latencia depende del escenario de propagación
o del nivel de QoS—, que es precisamente la estructura de dependencia
que la realidad problemática anticipa. El umbral del 85\% se introduce
aquí como \emph{criterio de éxito metodológico operacionalizable}:
habilita transformar un juicio cualitativo (“la red es de buena
calidad”) en una decisión verificable celda por celda dentro de la
matriz experimental.

{[}INSERTAR AQUÍ TABLA: Comparativa de los niveles de servicio MQTT (QoS
0/1/2) — confirmación, latencia, sobrecarga, confiabilidad y factor
experimental asociado{]}

\paragraph{Pregunta Específica 3
(PE3)}\label{pregunta-especuxedfica-3-pe3}

\textbf{¿Qué familia y configuración de modelo de Machine Learning
—implementada preferentemente con Scikit-learn sobre regresión lineal
regularizada, árboles de decisión y ensambles tipo \emph{random forest}—
establece la relación más robusta entre las variables físicas de
radiofrecuencia y las variables dinámicas de red, con qué exactitud
clasificaría el estado de QoS respecto del umbral del 85\% y qué
importancia relativa tendrían los predictores en la explicación de dicho
estado?}

\emph{Trasfondo de ingeniería.} PE3 se inscribe en el área de
\emph{aprendizaje automático aplicado a redes de comunicaciones}. El
problema puede formularse como aprendizaje supervisado de dos tipos: (i)
una tarea de \emph{clasificación}, cuya variable objetivo sería el
cumplimiento o incumplimiento del umbral del 85\% en las métricas de
QoS; y (ii) una tarea de \emph{regresión}, cuyo objetivo sería estimar
magnitudes continuas como la latencia o el jitter a partir de los
indicadores físicos y de configuración. La ingeniería de este problema
radica en que la relación entre RSSI y rendimiento no es monotónica ni
aditiva: conviven efectos de umbral (bajo cierto nivel de señal, la tasa
de error se desborda), efectos de saturación (por encima de un nivel
suficiente, mejorar la potencia no reduce la latencia) e interacciones
entre la carga y la calidad del enlace; los modelos lineales simples
capturan la tendencia general, mientras que los ensambles de árboles
capturan los efectos de umbral y las interacciones sin requerir
especificación previa de la forma funcional. La evaluación debe regirse
por buenas prácticas de validación: partición estratificada entre
entrenamiento y prueba, validación cruzada en \emph{k} pliegues,
búsqueda de hiperparámetros sin contaminar la partición de prueba, y
reporte conjunto de métricas (exactitud, precisión, sensibilidad,
\emph{F}1, área bajo la curva ROC, y MAE/RMSE/R² para la regresión) para
detectar tanto el sobreajuste como el sesgo por desequilibrio de clases.
Un requisito adicional, propio de la ingeniería aplicada, es la
\emph{interpretabilidad}: un modelo cuyos predictores no pueden leerse
en términos físicos (por ejemplo, que el RSSI o el nivel de QoS no
aparezcan entre los atributos más importantes) sería de escaso valor
para el gestor de red, aun cuando su exactitud sea elevada. PE3
convierte, así, el conjunto de mediciones en un \emph{modelo utilizable}
de correlación entre las dos familias de variables.

{[}INSERTAR AQUÍ FLUJOGRAMA DE: Esquema del aprendizaje supervisado
previsto — clasificación del umbral del 85\% y regresión de
latencia/jitter con Scikit-learn{]}

\paragraph{Pregunta Específica 4
(PE4)}\label{pregunta-especuxedfica-4-pe4}

\textbf{¿Es técnica y económicamente viable, y útil desde el punto de
vista de la operación, un prototipo IoT híbrido de bajo costo
—instrumentado con tableros de Power BI para la comunicación del estado
de QoS— que ofrezca una evaluación reproducible de la calidad de
servicio de la red interior, manteniendo el costo por nodo y la
complejidad de operación dentro de los rangos compatibles con
despliegues masivos y con la gestión cotidiana de instalaciones reales?}

\emph{Trasfondo de ingeniería.} PE4 cierra el sistema de preguntas con
la dimensión de \emph{ingeniería de sistemas y economía de las
telecomunicaciones}. La validación de un instrumento de medición no se
limita a demostrar que produce estimaciones consistentes en un banco de
pruebas, sino que exige una metodología reproducible y un valor de uso
demostrable: definición de escenarios de referencia (aulas, oficinas,
pasillos y salas técnicas), de KPIs comparables (PDR, latencia
percentil-95, jitter, buenuso, exactitud del modelo), de líneas base
técnicas (el mismo enlace sin instrumentación cruzada y las
arquitecturas centralizadas descritas en la realidad problemática) y de
un nivel de madurez tecnológica explícito. La simulación permite
explorar el espacio de diseño a escala, pero la ingeniería de
telecomunicaciones reconoce la brecha entre simulación y despliegue
real; por ello la pregunta incorpora el \emph{prototipo físico} como
etapa de verificación obligatoria. En paralelo, la viabilidad económica
se trata como variable de diseño: el CAPEX por nodo, el consumo de
energía, la ausencia de infraestructura centralizada y la dependencia de
componentes comerciales estándar deben cuantificarse para demostrar que
la solución se sitúa en la franja de costo exigida por el mercado IoT.
Finalmente, la dimensión de \emph{utilidad} se concreta en la capa de
visualización: los tableros de Power BI traducen las salidas del
análisis estadístico y del modelo de Machine Learning en indicadores
accionables —semáforos de cumplimiento del umbral, tendencias por
escenario y alertas de degradación—, de modo que el resultado no quede
confinado en el informe académico, sino que pueda sostener decisiones de
dimensionamiento, reubicación de puntos de acceso o ajuste de parámetros
del protocolo. En conjunto, PE4 convierte el resultado académico en una
\emph{propuesta de solución con grado de madurez tecnológica} explícito,
requisito indispensable de la ingeniería aplicada.

{[}INSERTAR AQUÍ IMAGEN TÉCNICA: Esquema de los tableros de Power BI —
tablero de estado con semáforos, tablero de diagnóstico y tablero
predictivo{]}

\subsubsection{Síntesis Integradora}\label{suxedntesis-integradora}

Las cuatro preguntas específicas no son compartimentos estancos, sino
eslabones de una cadena experimental coherente: PE1 (canal) suministra
la familia de variables físicas de radiofrecuencia; PE2 (protocolo y
red) produce la familia de variables dinámicas y define el umbral de
éxito dentro de una matriz factorial 3×3×3; PE3 (Machine Learning)
establece y valida la relación estadístico-aprendizaje entre ambas
familias; y PE4 (viabilidad) integra, instrumenta y somete a validación
técnica, económica y de uso el resultado sistémico. Responderlas en
conjunto permite, precisamente, dar respuesta a la pregunta general
sobre la capacidad de un prototipo IoT híbrido de bajo costo para
evaluar y comunicar de forma reproducible la calidad de servicio de las
redes inalámbricas interiores.

\subsection{1.3. Justificación}\label{justificaciuxf3n}

\subsubsection{1.3.1. Presentación del tema y del proceso
comunicativo}\label{presentaciuxf3n-del-tema-y-del-proceso-comunicativo}

La presente investigación propone el diseño, la implementación y la
evaluación de un prototipo IoT híbrido para la medición de la calidad de
servicio en redes inalámbricas interiores. El proceso comunicativo
objeto de estudio se circunscribe a una cadena específica: el nodo
sensor–actuador basado en ESP32 adquiere e incluye en cada publicación
tanto el estado medido como los indicadores de radiofrecuencia del
enlace; la trama viaja por el medio compartido conforme al protocolo
Wi-Fi IEEE 802.11 (MAC DCF/CSMA/CA y PHY 802.11 b/g/n) hasta el punto de
acceso; el \emph{broker} MQTT recibe, enruta y confirma las
publicaciones según el nivel de QoS contratado; el servidor de
adquisición registra las series temporales sincronizadas de variables
físicas y de red; y el canal de análisis —estadística descriptiva e
inferencial, modelos de Machine Learning y tableros de Power BI—
transforma esos registros en conocimiento sobre el desempeño del enlace.
Los actores involucrados son, por tanto, identificables: los
desarrolladores y fabricantes de dispositivos IoT (que requieren
criterios objetivos de desempeño), los administradores y gestores de
redes de área local (que necesitan diagnósticos accionables sobre
cobertura y congestión), la comunidad académica de telecomunicaciones
(que demanda protocolos experimentales replicables) y los usuarios
finales de instalaciones interiores (educativas, asistenciales,
industriales y comerciales), cuya experiencia de calidad depende
directamente del comportamiento del enlace.

{[}INSERTAR AQUÍ FLUJOGRAMA DE: Arquitectura general del sistema IoT
híbrido, mostrando la relación entre los nodos ESP32, el canal Wi-Fi, el
broker MQTT y la capa de análisis{]}

\subsubsection{1.3.2. Enfoques de la investigación y razones de su
elección}\label{enfoques-de-la-investigaciuxf3n-y-razones-de-su-elecciuxf3n}

La elección del tema responde a la convergencia de cuatro evidencias: el
crecimiento sostenido de la densidad de dispositivos IoT en interiores;
la concentración de las bandas ISM no licenciadas como espacio
compartido y congestionado; la brecha entre las soluciones de evaluación
de QoS de alto costo (analizadores de espectro, sondas comerciales,
plataformas de gestión empresarial) y la realidad económica de los
despliegues masivos; y la madurez suficiente de las técnicas de
aprendizaje automático de código abierto para extraer modelos
interpretables a partir de series temporales de sensores. Sobre esa
base, la investigación se aproxima a su objeto de estudio mediante los
siguientes enfoques:

\textbf{a) Enfoque teórico-científico.} Se busca producir conocimiento
verificable sobre la relación funcional entre variables físicas de
radiofrecuencia y variables dinámicas de red. En el plano conceptual, se
articulan marcos de referencia institucionales —la Recomendación ITU-T
E.800 para la definición de calidad de servicio, las Recomendaciones
ITU-R P.1238 e ITU-T G.114/Y.1540 para la propagación y el desempeño de
transferencia de paquetes, y el estándar IEEE 802.11 para el acceso al
medio— con el paradigma del aprendizaje automático supervisado. El
aporte teórico consiste en la formalización de un modelo de correlación
entre ambas familias de variables, comprobado empíricamente y
documentado con indicadores de ajuste e importancia de predictores, lo
que amplía el cuerpo de conocimiento sobre caracterización de QoS en
redes inalámbricas de interiores.

\textbf{b) Enfoque tecnológico.} La respuesta al problema se materializa
en un artefacto: un prototipo hardware–software de bajo costo,
replicable con componentes comerciales disponibles (microcontroladores
ESP32, punto de acceso Wi-Fi, \emph{broker} MQTT, base de datos y capa
de análisis). Este enfoque se elige porque la problemática descrita es
de naturaleza aplicada: la brecha identificada no se cierra con un
modelo analítico, sino con un instrumento que pueda instalarse, operarse
y repetirse en instalaciones reales. El prototipo incorpora desde su
diseño la instrumentación de las variables físicas, condición
indispensable para que la evaluación de QoS sea \emph{cross-layer} y no
una simple medición de rendimiento de aplicación.

\textbf{c) Enfoque metodológico.} Se opta por un diseño experimental de
carácter cuantitativo, estructurado como un plan factorial completo
3×3×3 con factores de propagación, de carga de tráfico y de nivel de
servicio MQTT, complementado con un plan de contingencia explícito de
cinco pasos para el tratamiento eventual de resultados que no alcancen
el umbral del 85\%. Este enfoque se justifica por tres razones: permite
estimar efectos principales e interacciones en una sola campaña de
medición; garantiza la reproducibilidad, al hacer públicos los niveles,
las réplicas y los criterios de éxito antes de la recolección; y evita
el sesgo de interpretación \emph{a posteriori}, al fijar por adelantado
tanto el criterio de éxito como el procedimiento de respuesta ante su
incumplimiento.

\textbf{d) Enfoque económico y social.} La investigación se delimita
explícitamente en la franja de costo de los despliegues IoT masivos: sin
infraestructura centralizada especializada, sin espectro licenciado y
con componentes de costo reducido. Su justificación social se deriva de
las aplicaciones que habilita: el dimensionamiento racional de redes en
aulas y laboratorios, la detección temprana de degradación en
instalaciones sanitarias y de seguridad, y la dotación de herramientas
de diagnóstico a instituciones que carecen de equipos de medición
especializados. Desde el punto de vista académico, se espera como aporte
la publicación de un protocolo experimental replicable y del conjunto de
datos estructurado resultante, disponible para trabajos posteriores de
validación y comparación.

\subsubsection{1.3.3. Justificación de la pertinencia del umbral del
85\% y de la
validación}\label{justificaciuxf3n-de-la-pertinencia-del-umbral-del-85-y-de-la-validaciuxf3n}

El umbral del 85\% se asume como criterio de éxito metodológico por su
carácter de frontera operacional verificable: permite pronunciarse, para
cada celda del diseño factorial y para el modelo predictivo, sobre si el
desempeño observado se considera aceptable, y habilita la activación de
un plan de contingencia de cinco pasos cuya finalidad no es su
redefinición retroactiva, sino la verificación del instrumental, el
re-aislamiento de causas, la ampliación de la evidencia, el reajuste de
parámetros configurables y, únicamente en último caso, la documentación
formal de la desviación. De este modo, la justificación metodológica del
estudio se sostiene tanto si el umbral se alcanza como si no se alcanza:
en el primer caso se dispone de evidencia de cumplimiento; en el
segundo, de evidencia de límite técnico correctamente diagnosticada.

\subsubsection{1.3.4. Delimitación del caso: coordenadas espaciales y
temporales}\label{delimitaciuxf3n-del-caso-coordenadas-espaciales-y-temporales}

En el plano espacial, la investigación se delimita a entornos interiores
de instalaciones de uso institucional —aulas, laboratorios, oficinas y
pasillos de circulación— caracterizados por geometrías cerradas,
mobiliario y tabiquería de materiales de construcción convencionales,
operando en la banda ISM de 2,4 GHz con infraestructura Wi-Fi de tipo
punto de acceso doméstico o institucional. Se descartan explícitamente
los escenarios exteriores, los de propagación entre plantas y las bandas
de 5 GHz y 6 GHz, así como los entornos con procesamiento de señales en
la capa física fuera del alcance del hardware empleado. En el plano
temporal, la investigación se delimita al periodo de ejecución previsto
en el cronograma del plan de tesis, organizado en cuatro etapas
consecutivas: (i) diseño y construcción del prototipo con pruebas de
integración; (ii) calibración del instrumento y campaña de mediciones
bajo el diseño factorial; (iii) procesamiento de datos, modelado con
Scikit-learn y construcción de tableros; y (iv) análisis, contrastación
de resultados, redacción y sustentación. Las mediciones se realizarán en
horarios y condiciones registradas, de modo que la campaña sea auditable
y susceptible de repetición posterior.

\subsection{1.4. Objetivos}\label{objetivos}

\subsubsection{1.4.1. Objetivo Principal}\label{objetivo-principal}

Diseñar, implementar y evaluar un prototipo IoT híbrido basado en
microcontroladores ESP32, enlace Wi-Fi IEEE 802.11 y publicaciones MQTT,
que permita correlacionar mediante técnicas de Machine Learning las
variables físicas de radiofrecuencia con las variables dinámicas de red,
con el fin de evaluar la calidad de servicio de las redes inalámbricas
interiores y verificar de manera reproducible el umbral de desempeño del
85\% fijado como criterio de éxito.

\subsubsection{1.4.2. Objetivos
Específicos}\label{objetivos-especuxedficos}

\textbf{OE1.} Caracterizar el canal de propagación en interiores de la
banda ISM de 2,4 GHz y las variables físicas de radiofrecuencia medibles
con el instrumento —RSSI, SNR, potencia de recepción y ruido de fondo—,
ajustando a las mediciones de campo un modelo de pérdida de distancia y
evaluando su bondad de ajuste.

\textbf{OE2.} Implementar el prototipo IoT híbrido —nodos ESP32, punto
de acceso Wi-Fi, \emph{broker} MQTT y estación de adquisición— con la
instrumentación y la sincronización necesarias para registrar de manera
simultánea y trazable las variables físicas y las variables dinámicas de
red.

\textbf{OE3.} Ejecutar la campaña de mediciones bajo el diseño factorial
experimental 3×3×3 y cuantificar las métricas de calidad de servicio
(latencia extremo a extremo, jitter, razón de entrega de tramas, pérdida
de paquetes y buen uso) frente al umbral del 85\%, aplicando el plan de
contingencia de cinco pasos ante resultados que no alcancen dicho
criterio.

\textbf{OE4.} Desarrollar y validar, con la biblioteca Scikit-learn, los
modelos de Machine Learning que correlacionen las variables físicas de
radiofrecuencia con las variables dinámicas de red, evaluando su
exactitud mediante validación cruzada, sus métricas de clasificación y
regresión, y la importancia relativa de sus predictores.

\textbf{OE5.} Construir los tableros de análisis en Power BI para la
visualización y comunicación del estado de calidad de servicio, y
documentar la viabilidad técnica y económica del prototipo frente a las
soluciones de infraestructura centralizada convencionales.

{[}INSERTAR AQUÍ TABLA: Matriz de consistencia de la investigación,
relacionando problema, objetivos, preguntas, hipótesis y variables{]}

{[}INSERTAR AQUÍ FLUJOGRAMA DE: Arquitectura general del sistema IoT
híbrido…{]}

\section{CAPÍTULO II: MARCO TEÓRICO Y
CONCEPTUAL}\label{capuxedtulo-ii-marco-teuxf3rico-y-conceptual}

El presente capítulo sustenta teóricamente la investigación. Se inicia
con los antecedentes de trabajos previos relacionados (2.1); continúa
con el desarrollo profundo de la base teórica que sustenta cada una de
las preguntas de investigación —calidad de servicio, Wi-Fi IEEE 802.11,
propagación en interiores, MQTT y Machine Learning— (2.2); procede al
estado del arte organizado por enfoques y cierra con el párrafo de vacío
que delimita la contribución pretendida (2.3); explicita el marco
conceptual o glosario rector (2.4) y finaliza con la normatividad
estrictamente pertinente al objeto de estudio (2.5).

\subsection{2.1. Antecedentes de la
Investigación}\label{antecedentes-de-la-investigaciuxf3n}

Los antecedentes se presentan diferenciados por enfoque, con el
propósito de mostrar el estado real del avance sobre el problema
principal y la distancia que separa a los trabajos existentes de la
propuesta formulada en este documento.

\subsubsection{2.1.1. Antecedentes sobre caracterización y modelado de
la propagación en
interiores}\label{antecedentes-sobre-caracterizaciuxf3n-y-modelado-de-la-propagaciuxf3n-en-interiores}

Chrysikos, Georgopoulos y Kotsopoulos (2009) realizaron una validación
específica del sitio del modelo de pérdida de la Unión Internacional de
Telecomunicaciones (ITU) a 2,4 GHz, a partir de mediciones adquiridas en
un entorno de oficina inalámbrica (\emph{WLAN}). Los autores lograron un
ajuste numérico del modelo ITU de carácter general al sitio particular,
obteniendo valores de los parámetros del modelo correspondientes a esa
banda de frecuencia, con el propósito de mejorar la precisión del
presupuesto de enlace en redes inalámbricas locales. Su aporte evidencia
que los modelos de referencia públicos requieren recalibración local
para ser utilizados como herramienta de diseño.

Solahuddin y Mardeni (2011) desarrollaron un modelo empírico de pérdida
para redes IEEE 802.11n a 2,4 GHz en un edificio académico, comparando
el desempeño de los modelos de doble pendiente, segmentado, de sombreado
log-normal, ITU-R P.1238-1, Motley-Keenan ajustado y COST 231 de
múltiples paredes. Concluyeron que el modelo segmentado ofrecía la mejor
estimación de señal y lo tomaron como referencia para una optimización
posterior. Mardeni y Solahuddin (2012) continuaron ese trabajo,
reportando que el modelo mejorado reducía el error medio de predicción
respecto del modelo original y que la validación en un edificio de
oficinas distinto mantenía la precisión obtenida.

Kaddouri, El Hajj, Zaharia y El Zein (2018) efectuaron mediciones
extensas de pérdida de trayecto en una oficina de planta abierta a 2,4
GHz y 5,8 GHz, cubriendo combinaciones de posición y altura
transmisor–receptor en visibilidad directa y sin visibilidad directa, y
realizaron una segunda campaña para cuantificar el efecto de la
movilidad de las personas. Reportaron que la presencia de las personas
incrementa la pérdida media del trayecto y que los exponentes de pérdida
obtenidos difieren según la configuración de medición, además de
contrastar sus resultados con los valores de la recomendación ITU-R
P.1238.

Rath, Timmadasari, Panigrahi y Simha (2017) propusieron un modelo
estadístico de pérdida en interiores para operaciones regulares de Wi-Fi
en la banda de 2,4 a 2,5 GHz, validado mediante campañas de medición en
oficinas con distintas condiciones (espacio ocupado, oficina abierta,
pasillo, comedor y locales de varios pisos). Compararon su modelo con el
de ITU-R y con el de distancia logarítmica, reportando mejor ajuste de
su propuesta a los datos experimentales y destacando que los exponentes
de pérdida varían de forma significativa entre pasillos y espacios
abiertos debido a la multipropagación y a las reflexiones.

{[}INSERTAR AQUÍ TABLA: Cuadro comparativo de antecedentes de
propagación en interiores — autores, banda, modelo evaluado y hallazgo
clave{]}

\subsubsection{2.1.2. Antecedentes sobre medición experimental del
desempeño Wi-Fi y del QoS de protocolos de
aplicación}\label{antecedentes-sobre-mediciuxf3n-experimental-del-desempeuxf1o-wi-fi-y-del-qos-de-protocolos-de-aplicaciuxf3n}

Dianu, Riihijärvi y Petrova (2014) presentaron un estudio basado en
mediciones del desempeño del estándar Wi-Fi IEEE 802.11ac en un entorno
interior típico de oficina. Observaron que las mejoras de desempeño
respecto de 802.11n son sensibles a las condiciones del canal —con caída
rápida de las tasas de datos al aumentar la distancia— y que, mientras
la interferencia de canal adyacente proveniente de equipos legados tiene
un impacto severo, el mecanismo de acceso al medio permite compartir el
canal de manera efectiva cuando la interferencia es de canal
coexistente.

Fiehe, Riihijärvi y Mähönen (2010) llevaron a cabo campañas de medición
extensas de IEEE 802.11n en entornos de oficina controlados e
interferidos. Reportaron que el rendimiento teórico no se alcanza en la
práctica, que 802.11n tolera de manera razonable la interferencia
proveniente de otras tecnologías de la banda ISM de 2,4 GHz, pero que
interferentes estrechos de potencia suficiente y equipos legados en la
misma región de espectro degradan de manera notable el desempeño.

Miaoudakis, Stratakis, Antonidakis, Zaharopoulos y Stojanovic (2009)
evaluaron mediante mediciones reales la coexistencia de redes
inalámbricas locales y de área personal en una oficina típica,
considerando interferencia 802.11–802.11 en canales solapados y no
solapados y la interferencia Bluetooth sobre 802.11. Reportaron
degradaciones importantes de rendimiento y aumentos de la latencia de
ida y vuelta cuando los canales se solapan, y degradación residual aun
en canales no solapados.

van Bloem, Schiphorst, Kluwer y Slump (2012) desarrollaron una
herramienta de medición de la influencia de distintas fuentes de
interferencia sobre enlaces de comunicación IEEE 802.11, analizando la
capa física y la capa de enlace para evaluar el impacto conjunto sobre
la calidad de servicio. Sus resultados mostraron que ciertos
interferentes producen degradación severa del QoS del Wi-Fi, y que el
análisis a nivel de tramas permite identificar el mecanismo de
interferencia con mayor detalle que las herramientas tradicionales
centradas únicamente en rendimiento y tasa de error de paquetes.

Parastar y Wang (2020) estudiaron experimentalmente la calidad de
servicio en redes IEEE 802.11, poniendo énfasis en el mecanismo de
acceso mejorado EDCA propuesto por IEEE 802.11e y en la influencia de
sus parámetros (AIFS y TXOP) sobre el retardo, empleando para ello una
plataforma real con redes definidas por software, en contraposición a
los estudios mayoritariamente basados en simulación.

En el plano de los protocolos de aplicación, Hijazi, Habaebi, Al-Haddad
y Mohammed Zyoud realizaron una prueba de estrés de un servidor MQTT en
una infraestructura de red privada de IoT, analizando retardo, jitter,
pérdida de paquetes y rendimiento en escenarios de tiempo real y no
tiempo real para los tres niveles de servicio MQTT. Reportaron que el
retardo crece de forma marcada con la carga, que el nivel QoS 0 presenta
los menores retardos por ausencia de confirmaciones y que los niveles
QoS 1 y QoS 2 intercambian sobrecarga de control por mayor confiabilidad
de entrega.

Akbar, Amron, Mulya y Hanifah (2017) revisaron la implementación del
protocolo MQTT para la comunicación en redes de sensores inalámbricas,
midiendo el proceso de descubrimiento de puerta de enlace, la pérdida de
paquetes en QoS 0, la duplicación de datos y los reintentos en QoS 1 y
2, y el tiempo de ida y vuelta en los niveles confirmados, concluyendo
que el nivel de servicio elegido determina explícitamente el compromiso
entre velocidad y confiabilidad de la entrega.

Nikfarjam Nashtifani (2026) midió experimentalmente la latencia en
sistemas IoT de pequeña escala construidos con nodos ESP32 que publican
por MQTT sobre Wi-Fi en interiores, sincronizando los relojes mediante
NTP para obtener retardos de ida y vuelta de un solo sentido y de ida y
vuelta, en distintas distancias de comunicación y en los niveles de QoS
0 y 1. Concluyó que las condiciones de la red Wi-Fi y la arquitectura de
comunicación influyen en la latencia en mayor medida que las
limitaciones de proceso del microcontrolador, y que el nivel QoS 1
incrementa el retardo y el jitter respecto de QoS 0.

Chen y Kunz, por su parte, compararon experimentalmente protocolos de
IoT bajo una red inalámbrica de acceso con restricciones de ancho de
banda, retardo y pérdida, evidenciando que el comportamiento de MQTT
ante la pérdida de paquetes y la reducción de capacidad del enlace lo
distancia del desempeño de alternativas orientadas a la mensajería de
baja latencia.

{[}INSERTAR AQUÍ TABLA: Cuadro comparativo de antecedentes de medición
Wi-Fi y QoS MQTT — autores, protocolo, métricas y hallazgo{]}

\subsubsection{2.1.3. Antecedentes sobre Machine Learning aplicado a la
predicción de QoS y al estado de los
enlaces}\label{antecedentes-sobre-machine-learning-aplicado-a-la-predicciuxf3n-de-qos-y-al-estado-de-los-enlaces}

Pawar, Jadhav, Lokhande, Raskar y Patil (2023) evaluaron los parámetros
de calidad de servicio de la comunicación MQTT en aplicaciones de IoT
mediante una red neuronal profunda, con el propósito de anticipar el
comportamiento del enlace a partir de series de indicadores de red.

Hameed (2023) planteó un marco de modelado de tráfico y predicción de
QoS mediante técnicas de aprendizaje profundo avanzado, formulando la
predicción de indicadores como problema de pronóstico de series
temporales univariadas y multivariadas y validando la propuesta con
datos de tráfico de aplicaciones IoT reales sobre acceso IEEE 802.15.4.

Çoğay, Akkoç y Seçinti (2024) propusieron un marco de enrutamiento
dinámico multiproto asistido por Machine Learning en el que un
clasificador de paquetes basado en métodos de ensamble selecciona las
rutas de acuerdo con la carga y el estado de los enlaces, evaluando la
clasificación con métricas F1 y el enrutamiento en el simulador NS-3
mediante razones de entrega de paquetes y rendimiento.

Ávalos-Arce, Pérez-Díaz, Del-Valle-Soto y Briseño (2024) investigaron
las limitaciones de los modelos de predicción del estado de paquetes en
redes inalámbricas IEEE 802.15.4 sometidas a distintas interferencias,
entrenando varios clasificadores supervisados con características
extraídas de la captura de tramas. Reportaron que los atributos de
fortaleza de red y de tamaño de paquete resultan predictores dominantes
y que los métodos de ensamble (refuerzo de gradiente y bosques
aleatorios) obtienen el mejor desempeño frente a métodos lineales y
basados en vectores de soporte.

Yazici y Güreş (2024) predijeron indicadores de QoS —razón de entrega de
paquetes y rendimiento— en escenarios inalámbricos empleando múltiples
métodos de Machine Learning con un esquema de validación cruzada anidada
que evita la fuga de información entre la selección de parámetros y la
evaluación; encontraron que los métodos de ensamble, y en particular los
bosques aleatorios, ofrecen el mejor compromiso de error y coeficiente
de determinación.

Becvar, Plachy, Mach, Nikolov y Gesbert (2024) abordaron la predicción
de la calidad de canal con Machine Learning, desde el concepto hasta la
validación experimental, en un trabajo que subraya la brecha entre el
desempeño reportado en condiciones controladas y el alcanzado en
condiciones operativas reales.

Benadji, Zitoune y Vèque (2023) propusieron la predicción de la razón de
pérdida de paquetes como indicador de congestión en redes IoT mediante
modelos de aprendizaje profundo y ventanas deslizantes sobre series
temporales, con el fin de sustituir la reacción tardía a la congestión
por una anticipación basada en datos.

{[}INSERTAR AQUÍ TABLA: Cuadro comparativo de antecedentes de Machine
Learning para predicción de QoS — autores, método, predictores y
métricas{]}

\subsubsection{2.1.4. Antecedentes sobre monitoreo, visualización y
tableros de
indicadores}\label{antecedentes-sobre-monitoreo-visualizaciuxf3n-y-tableros-de-indicadores}

Ehrlich, Neumann, Biendarra y Jasperneite (2019) abordaron el monitoreo
de la calidad de servicio en redes de comunicación industriales
híbridas, analizando los métodos disponibles de extracción de
información del tráfico y su aprovechamiento por los sistemas de gestión
de red, en un contexto donde la visibilidad del estado de la red es
condición para la toma de decisiones de configuración.

Calderon, del Campo, Saavedra y Santamaría (2023) implementaron y
evaluaron el desempeño de una plataforma IoT modular de arquitectura
distribuida que combina herramientas de la familia Elastic con Apache
Kafka, incluyendo un marco de monitoreo basado en agentes que supervisa
la plataforma con arreglo a métricas definidas y la complementa con
interfaces de visualización para facilitar el despliegue, la gestión y
el mantenimiento.

Praja, Tahir, Cahyani, Zahroeva, Anam y Prameswari (2025) desarrollaron
un sistema de monitoreo y registro del tráfico Wi-Fi de aulas
universitarias mediante captura de paquetes con Wireshark y su
visualización interactiva en Power BI, identificando patrones de uso,
tráfico dominante y anomalías potenciales a partir de los datos
capturados.

\subsubsection{2.1.5. Síntesis de los
antecedentes}\label{suxedntesis-de-los-antecedentes}

Los antecedentes descritos cubren, de manera conjunta, los cuatro
frentes que la investigación moviliza: la caracterización física del
canal, la medición experimental del desempeño Wi-Fi y de los protocolos
de aplicación, la predicción del comportamiento del enlace mediante
aprendizaje automático y la visualización de indicadores para la toma de
decisiones. Sin embargo, cada conjunto de trabajos desarrolla su aporte
de manera aislada respecto de los otros: los trabajos de propagación
entregan modelos de atenuación sin vincularlos con métricas de QoS de
aplicación; los de medición de protocolos entregan métricas sin modelar
su dependencia de las variables físicas; los de Machine Learning
predicen indicadores sin instrumentar de forma simultánea el plano
físico con hardware de bajo costo; y los de visualización presentan
tableros sin corresponder a modelos predictivos validados
experimentalmente.

{[}INSERTAR AQUÍ TABLA: Matriz de vacío de investigación — cobertura de
los cuatro frentes en la literatura frente a la propuesta{]}

\subsection{2.2. Marco Teórico}\label{marco-teuxf3rico}

\subsubsection{2.2.1. Calidad de Servicio
(QoS)}\label{calidad-de-servicio-qos}

La Recomendación ITU-T E.800 define la calidad de servicio como el
\emph{conjunto de efectos de un servicio que condicionan la satisfacción
del usuario del mismo}. La definición, deliberadamente amplia, tiene una
consecuencia práctica decisiva para la ingeniería: para que la QoS sea
gestionable debe traducirse en \emph{indicadores medibles} (KPI)
relacionados con la percepción del usuario o con el cumplimiento de
acuerdos de nivel de servicio. De esta manera, la QoS opera en tres
planos diferenciados: el plano del usuario (percepción subjetiva), el
plano del servicio (parámetros contractuales) y el plano de la red
(magnitudes técnicas observables en los equipos). La investigación se
sitúa en el tercer plano, como fuente de evidencia que sustenta los dos
anteriores.

Las magnitudes técnicas empleadas en este trabajo se definen con
precisión operacional:

\begin{itemize}
\tightlist
\item
  \textbf{Latencia extremo a extremo (\emph{end-to-end delay}).} Tiempo
  transcurrido entre la emisión de un mensaje por el nodo publicador y
  su recepción efectiva por el suscriptor, incluyendo el tiempo de
  acceso al medio Wi-Fi, la cola de espera en el \emph{broker} y el
  procesamiento de los extremos. La Recomendación ITU-T G.114 establece
  que un retardo de ida de hasta 150 ms es prácticamente imperceptible
  para la mayoría de los usuarios en aplicaciones interactivas, mientras
  que por encima de 400 ms resulta generalmente inaceptable; este rango
  se utiliza en la investigación como referencia interpretativa, no como
  umbral de éxito.
\item
  \textbf{Jitter o variación del retardo.} Dispersión estadística de los
  retardos de entrega consecutivos, típicamente expresada como la
  desviación típica de la serie de retardos o como la diferencia de
  retardo de paquetes (IPDV, definido en la RFC 3393). Un jitter elevado
  degrada las aplicaciones basadas en flujo continuo aun cuando la
  latencia media sea aceptable.
\item
  \textbf{Razón de entrega de tramas (PDR, \emph{Packet Delivery
  Ratio}).} Cociente entre el número de mensajes recibidos correctamente
  y el número de mensajes publicados, expresado en porcentaje. Es el
  indicador binario de confiabilidad de extremo a extremo y el natural
  para expresar el umbral del 85\% adoptado en la investigación.
\item
  \textbf{Pérdida de paquetes.} Complemento probabilístico de la PDR,
  asociada a colisiones en el medio, a desbordamiento de colas, a
  degradación del enlace físico y a desconexiones temporales.
\item
  \textbf{Buen uso (\emph{goodput}).} Cantidad de carga útil entregada
  por unidad de tiempo, excluyendo las cabeceras, las confirmaciones y
  los mensajes retransmitidos; permite distinguir el rendimiento
  aparente del rendimiento efectivamente aprovechable por la aplicación.
\item
  \textbf{Índices de radiofrecuencia.} RSSI (\emph{Received Signal
  Strength Indicator}), SNR (\emph{Signal-to-Noise Ratio}), nivel de
  ruido de fondo y, cuando el controlador lo expone, la tasa de error
  estimada. Constituyen la familia de \emph{variables físicas} del
  modelo de correlación.
\end{itemize}

Conviene subrayar que la QoS es una propiedad \emph{emergente} del
sistema completo: no se explica por la calidad del punto de acceso, ni
por el ancho de banda contratado, ni por el protocolo de aplicación
considerados aisladamente, sino por la interacción entre el canal de
propagación, el mecanismo de acceso al medio y la semántica de
confirmación del protocolo de transporte de mensajes. Esta
característica justifica formalmente el enfoque transversal de la
investigación.

{[}INSERTAR AQUÍ TABLA: Definición operacional de las métricas de QoS
empleadas — latencia, jitter, PDR, pérdida y buen uso, con unidad y
referencia normativa{]}

\subsubsection{2.2.2. Wi-Fi IEEE 802.11: acceso al medio y parámetros de
operación}\label{wi-fi-ieee-802.11-acceso-al-medio-y-paruxe1metros-de-operaciuxf3n}

La familia de estándares IEEE 802.11 define las capas física y de enlace
para redes inalámbricas locales (WLAN). En la banda ISM de 2,4 GHz
operan las variantes 802.11b, 802.11g y 802.11n, esta última con
modulaciones OFDM, canales de 20 y 40 MHz y opciones de espacial
\emph{Multiple-Input Multiple-Output}; las variantes 802.11ac y 802.11ax
operan predominantemente en la banda de 5 GHz y, en el caso de 802.11ax,
también en 6 GHz. Los canales de 20 MHz en 2,4 GHz se solapan entre sí,
de modo que únicamente tres canales son mutuamente no solapados
(típicamente los canales 1, 6 y 11), condición que constituye una
restricción de diseño de primer orden en instalaciones con múltiples
puntos de acceso y con presencia de redes vecinas.

En la capa de enlace, el mecanismo básico de acceso distribuido (DCF,
\emph{Distributed Coordination Function}) implementa un procedimiento de
acceso al medio por detección de portadora con evitación de colisiones
(CSMA/CA). Antes de transmitir, la estación verifica que el medio esté
inactivo durante un intervalo de escucha extendida; si el medio está
ocupado o si se produce una colisión, difiere la transmisión mediante un
\emph{retroceso exponencial binario}: la estación selecciona
aleatoriamente un entero dentro de una ventana de contienda cuyo tamaño
se duplica hasta alcanzar un límite máximo, y vuelve a intentar el
acceso. Este mecanismo se modela analíticamente mediante cadenas de
Markov de dos dimensiones, en la formulación clásica de Bianchi (2000),
que permite derivar expresiones de la probabilidad de transmisión
condicionada, la probabilidad de colisión y la tasa de éxito de tramas
en función de los parámetros de la ventana de contienda y del número de
estaciones activas.

Tres consecuencias de este diseño son directamente relevantes para la
investigación:

\begin{enumerate}
\def\labelenumi{\arabic{enumi}.}
\tightlist
\item
  \textbf{Sensibilidad a la densidad y a la carga.} A medida que el
  número de emisores activos o la tasa de ofrecimiento aumentan, la
  probabilidad de colisión crece de manera no lineal y la ventana de
  contienda se expande, elevando la latencia y reduciendo el buen uso.
  La congestión, por tanto, no es únicamente un fenómeno de red, sino un
  fenómeno del propio mecanismo de acceso.
\item
  \textbf{Sensibilidad a la calidad del canal.} Los mecanismos de sondeo
  de portadora, la selección de la tasa de modulación y codificación y
  las retransmisiones enlace a enlace reaccionan al RSSI y al SNR; una
  degradación física se traduce en aumento de la tasa de error, mayor
  número de reintentos y, en consecuencia, mayor jitter y menor PDR. Es
  aquí donde las variables físicas y las variables de red se acoplan.
\item
  \textbf{Sobrecarga y asimetría de interferencias.} La banda de 2,4 GHz
  es compartida con Bluetooth, Zigbee, cámaras inalámbricas, enlaces de
  audio y video y equipos vecinos no coordinados. Las interferencias de
  canal adyacente suelen ser más dañinas que las de canal coexistente,
  por cuanto el mecanismo de acceso puede coordinarse únicamente con las
  estaciones que perciben el mismo canal.
\end{enumerate}

Para dar respuesta a las aplicaciones con requerimientos diferenciados,
IEEE 802.11e introdujo el mecanismo EDCA (\emph{Enhanced Distributed
Channel Access}), comercializado como WMM, que sustituye una única cola
de contienda por cuatro categorías de acceso con parámetros de escucha,
ventana de retroceso y duración de transmisión diferenciados,
permitiendo priorizar el tráfico de voz, vídeo, \emph{best effort} y de
fondo. EDCA no garantiza QoS, sino que la \emph{diferencia}: la garantía
efectiva sigue dependiendo del canal y de la carga, lo que mantiene
vigente la necesidad de medición experimental.

{[}INSERTAR AQUÍ FLUJOGRAMA DE: Mecanismo DCF/CSMA/CA con retroceso
exponencial binario y su interacción con la calidad del canal
(cross-layer){]}

\subsubsection{2.2.3. Propagación en
interiores}\label{propagaciuxf3n-en-interiores}

El canal inalámbrico interior se caracteriza por tres fenómenos
dominantes: la \emph{pérdida de trayecto} asociada a la expansión de la
onda, la \emph{atenuación por obstáculos} —muros, tabiques, mobiliario,
conductos y cuerpos humanos— y el \emph{desvanecimiento por
multipropagación}, causado por la interferencia entre las múltiples
copias de la señal que llegan al receptor por reflexión, difracción y
dispersión. Cuando el receptor recibe una componente dominante
acompañada de un conjunto de componentes dispersas, el desvanecimiento
se describe con distribución Rician parametrizada por el factor K; en
ausencia de componente dominante, corresponde a la distribución de
Rayleigh; la familia Nakagami-\emph{m} generaliza ambos casos. Sobre la
envolvente de potencia media se superpone el \emph{desvanecimiento por
sombra}, de distribución log-normal, que captura la variación lenta de
la atenuación por el desplazamiento respecto de los obstáculos.

El modelo de pérdida de trayecto de distancia logarítmica constituye la
herramienta de trabajo por excelencia de la investigación:

PL(d) = PL(d₀) + 10 · n · log₁₀(d / d₀) + Xσ

donde PL(d) es la pérdida de trayecto en decibelios a la distancia d, d₀
es una distancia de referencia, n es el exponente de pérdida de trayecto
y Xσ es una variable aleatoria gaussiana de desviación típica σ que
representa el desvanecimiento por sombra. La Recomendación ITU-R P.1238
proporciona una formulación específica para interiores de la forma L =
20 log₁₀(f) + N log₁₀(d) + Lf(n) − 28 dB, donde f es la frecuencia en
megahercios, d la separación en metros, N el coeficiente de pérdida de
distancia y Lf(n) la pérdida por penetración de plantas, con valores
tabulados de N para contextos residenciales, de oficina y comercial. Los
modelos de múltiples paredes, en la línea de Motley-Keenan, añaden un
término sumatorio que contabiliza el número y el tipo de materiales
atraviesados por el trayecto.

Dos cuestiones de ingeniería condicionan el aprovechamiento de estos
modelos en la investigación:

\begin{itemize}
\tightlist
\item
  \textbf{Calibración local.} Los coeficientes N, σ y las pérdidas por
  material son específicos del edificio y de la banda. Los antecedentes
  revisados coinciden en que el modelo general debe ajustarse a las
  mediciones del sitio antes de emplearse como herramienta de
  predicción, y que el ajuste reduce el error de predicción de forma
  apreciable.
\item
  \textbf{Adecuación al instrumento.} El prototipo basado en ESP32 no
  dispone de instrumentación de canal complejo; expone fundamentalmente
  el RSSI del controlador Wi-Fi, a partir del cual es posible estimar el
  nivel de potencia recibida y el ruido de fondo. Por ello, la
  caracterización se plantea en términos de magnitudes
  \emph{observables} con el hardware disponible, y la fidelidad del
  modelo se evalúa por su capacidad de reproducir la tendencia y la
  dispersión de esas observaciones, no por su capacidad de reconstruir
  el campo electromagnético completo.
\end{itemize}

La multipropagación posee además una consecuencia no lineal que
justifica el enfoque de aprendizaje automático: el RSSI no disminuye
monotónicamente con la distancia en escenarios cerrados, ya que las
interferencias constructivas y destructivas producen máximos y mínimos
locales. Un modelo lineal simple captura la tendencia global, pero no el
comportamiento local que condiciona la calidad del enlace a distancias
cortas.

{[}INSERTAR AQUÍ TABLA: Cuadro comparativo de modelos de propagación en
interiores — ITU-R P.1238 vs Motley-Keenan vs log-distance vs modelos
estocásticos (Rayleigh, Rician, Nakagami-m){]}

\subsubsection{2.2.4. MQTT y la arquitectura de transporte de mensajes
IoT}\label{mqtt-y-la-arquitectura-de-transporte-de-mensajes-iot}

MQTT (\emph{Message Queuing Telemetry Transport}) es un protocolo de
mensajería ligero basado en el paradigma publicación–suscripción,
definido como estándar por OASIS (versión 3.1.1 en 2014 y versión 5.0 en
2019) y concebido para entornos de recursos limitados y enlaces de baja
fiabilidad. Su arquitectura se organiza en tres entidades: el
\emph{cliente publicador}, que envía mensajes a un \emph{tema}
(\emph{topic}) jerárquico; el \emph{cliente suscriptor}, que se suscribe
a temas de interés; y el \emph{broker} (servidor de mensajería), que
mantiene las sesiones, enruta los mensajes y gestiona las
confirmaciones. La separación entre publicadores y suscriptores elimina
la necesidad de que los extremos se conozcan mutuamente, lo que
simplifica el despliegue de nodos IoT y habilita la incorporación de
nuevos consumidores de datos —como la base de datos de la investigación
y la capa de tableros— sin modificar la configuración de los emisores.

El protocolo define tres niveles de servicio que constituyen el eje de
la variable experimental “nivel de QoS”:

\begin{itemize}
\tightlist
\item
  \textbf{QoS 0 (al menos una vez, \emph{at most once}).} Entrega sin
  confirmación: el mensaje se envía una sola vez y no se retransmite. Es
  el nivel más rápido y de menor sobrecarga, pero no garantiza la
  recepción ni la no duplicación.
\item
  \textbf{QoS 1 (al menos una vez).} Entrega con confirmación de
  recibido (PUBACK): el emisor retransmite hasta recibir la
  confirmación, lo que garantiza la entrega al menos una vez, a costa de
  mayor latencia, mayor consumo de ancho de banda y posibilidad de
  duplicados.
\item
  \textbf{QoS 2 (exactamente una vez).} Ciclo completo de cuatro
  mensajes (PUBREC/PUBREL/PUBCOMP) que evita la duplicación mediante
  control explícito de identificadores de paquete. Es el nivel con mayor
  sobrecarga de control y el mayor tiempo de entrega.
\end{itemize}

A los niveles de servicio se suman parámetros de sesión con impacto
directo en las métricas: el intervalo de \emph{keep alive} (que
determina la detección de desconexión), el indicador de sesión limpia
(que decide la recuperación de mensajes pendientes), el tiempo de
caducidad de sesión, el tamaño máximo de carga útil y el uso de marcas
de retención (\emph{retain}) y de mensajes de último valor conocido
(\emph{last will}). La versión 5.0 añade \emph{reason codes},
\emph{acknowledgements} diferidos y propiedades de mensaje que permiten
afinar el control de flujo y la depuración, características que se
consideran en la elección del \emph{broker}.

La elección de MQTT como protocolo de la investigación obedece a tres
razones de ingeniería: es el protocolo de facto del IoT sobre IP,
dispone de implementaciones de \emph{broker} ligeras y de código abierto
utilizables en una estación de adquisición local, y expone de forma
explícita y configurable el compromiso entre confiabilidad y sobrecaba,
lo que lo convierte en un factor experimental manipulable directamente.
Sobre Wi-Fi, MQTT opera sobre TCP, de modo que la latencia extremo a
extremo acumula el tiempo de acceso al medio inalámbrico, la ventana de
contienda, las retransmisiones TCP y el ciclo de confirmación MQTT: el
valor medido es, por construcción, una magnitud \emph{transversal} que
integra la pila completa.

El ESP32 constituye la plataforma de nodo elegida por su integración
nativa de interfaz Wi-Fi IEEE 802.11 b/g/n y Bluetooth, su capacidad de
operar como estación y como punto de acceso, su costo reducido, la
disponibilidad de bibliotecas MQTT maduras y su perfil de consumo
compatible con alimentación por bus o por batería de mediana duración.
El prototipo emplea el ESP32 tanto como nodo publicador como
instrumentador de los índices de radiofrecuencia, aprovechando la
interfaz de estado del controlador Wi-Fi para obtener RSSI y nivel de
ruido en cada publicación.

{[}INSERTAR AQUÍ FLUJOGRAMA DE: Modelo Publish-Subscribe de MQTT con
flujos de confirmación de los niveles QoS 0, 1 y 2{]}

\subsubsection{2.2.5. Machine Learning aplicado a la evaluación de
redes}\label{machine-learning-aplicado-a-la-evaluaciuxf3n-de-redes}

El aprendizaje automático comprende los métodos que permiten a un
sistema mejorar su desempeño a partir de la experiencia, en oposición a
la programación explícita de reglas. La investigación se circunscribe al
\emph{aprendizaje supervisado}, en el que se dispone de un conjunto de
datos de entrada \textbf{X} (predictores) y de una variable objetivo
\textbf{y} (respuesta) y se busca la función f : \textbf{X} → \textbf{y}
que generalice a datos no observados. En el problema planteado se
formulan dos tareas complementarias:

\begin{itemize}
\tightlist
\item
  \textbf{Clasificación.} La variable objetivo es discreta: por ejemplo,
  si la combinación experimental cumple o no cumple el umbral del 85\%
  en la métrica de referencia (PDR), o el estado de QoS codificado como
  adecuado / degradado / crítico.
\item
  \textbf{Regresión.} La variable objetivo es continua: la predicción de
  la latencia, del jitter o del buen uso a partir de los índices de
  radiofrecuencia y de las condiciones de operación.
\end{itemize}

Las familias de algoritmos consideradas se seleccionan por su idoneidad
con conjuntos de datos tabulares de tamaño moderado y por su
interpretabilidad:

\begin{itemize}
\tightlist
\item
  \textbf{Regresión lineal con regularización (Ridge y Lasso).}
  Proporcionan una línea base de referencia, cuantifican la asociación
  lineal entre cada predictor y la respuesta y, mediante la
  penalización, permiten reducir la varianza de las estimaciones e
  identificar predictores poco informativos.
\item
  \textbf{Árbol de decisión.} Modelan relaciones no lineales e
  interacciones por particiones sucesivas del espacio de predictores; su
  lectura es directa, pero presentan varianza alta y tendencia al
  sobreajuste cuando se profundiza el árbol.
\item
  \textbf{Bosques aleatorios y otros ensambles (\emph{gradient
  boosting}).} Combina múltiples árboles entrenados sobre muestras y
  subconjuntos de atributos, reduciendo la varianza sin incrementar
  proporcionalmente el sesgo; se reportan ampliamente como los de mejor
  desempeño en problemas de predicción de estado de enlaces e
  indicadores de red.
\item
  \textbf{Métodos lineales logísticos para clasificación.} Actúan como
  modelo de referencia interpretable ante problemas binarios.
\end{itemize}

La evaluación de los modelos se guía por principios que la investigación
adopta de forma explícita:

\begin{enumerate}
\def\labelenumi{\arabic{enumi}.}
\tightlist
\item
  \textbf{Separación entre entrenamiento y prueba.} Los datos se
  particionan, con estratificación cuando las clases están
  desequilibradas, de modo que la evaluación final se realice sobre
  observaciones no vistas durante el ajuste.
\item
  \textbf{Validación cruzada en \emph{k} pliegues.} Permite estimar la
  varianza del desempeño y detectar inestabilidad del modelo; en
  problemas con búsqueda de hiperparámetros se prefiere el esquema
  anidado, que evita la \emph{fuga} de información entre la selección de
  configuraciones y la evaluación.
\item
  \textbf{Métricas acordes al desbalance de clases.} Además de la
  exactitud, se emplean precisión, sensibilidad, especificidad,
  \emph{F}1 y el área bajo la curva ROC; para la regresión, el error
  absoluto medio, la raíz del error cuadrático medio y el coeficiente de
  determinación.
\item
  \textbf{Interpretabilidad.} El análisis de importancia de atributos y
  las correlaciones parciales se emplean para traducir el modelo al
  lenguaje físico de la ingeniería: si el RSSI, la carga de tráfico y el
  nivel de QoS no aparecen entre los predictores dominantes, el modelo
  será cuestionado aun con métricas elevadas.
\item
  \textbf{Reproducibilidad.} La fijación de semillas aleatorias, la
  versionación del conjunto de datos y la conservación del
  \emph{pipeline} completo de preprocesamiento y modelado forman parte
  del protocolo experimental.
\end{enumerate}

La biblioteca Scikit-learn (Pedregosa y otros, 2011) se elige como
entorno de implementación por su madurez, su disponibilidad de código
abierto, su coherencia de API entre preprocesamiento, selección de
modelos y métricas, y su compatibilidad con el ecosistema de análisis de
datos empleado en la investigación.

{[}INSERTAR AQUÍ FLUJOGRAMA DE: Pipeline de evaluación de modelos de
Machine Learning — partición, validación cruzada en k pliegues, métricas
e interpretabilidad{]}

\subsection{2.3. Estado del Arte}\label{estado-del-arte}

El estado del arte se organiza por enfoques, de modo que cada frente de
trabajo aparezca con su alcance real, sus resultados consolidados y sus
límites.

\subsubsection{Enfoque 1: Modelado y calibración de la propagación en
interiores}\label{enfoque-1-modelado-y-calibraciuxf3n-de-la-propagaciuxf3n-en-interiores}

Las campañas de medición realizadas en oficinas, edificios académicos y
locales comerciales a 2,4 GHz han consolidado un conjunto de
conocimientos estables: el exponente de pérdida de trayecto varía
significativamente entre espacios abiertos, pasillos y ambientes con
tabiquería; la presencia de personas incrementa de forma medible la
atenuación media; los coeficientes de las recomendaciones generales
deben ajustarse al sitio para reducir el error de predicción; y los
modelos de múltiples paredes superan a los modelos de distancia simple
cuando el trayecto atraviesa obstáculos identificables (Chrysikos y
otros, 2009; Solahuddin y Mardeni, 2011; Mardeni y Solahuddin, 2012;
Kaddouri y otros, 2018; Rath y otros, 2017). Este enfoque alcanza, por
tanto, un grado de madurez alto en lo que respecta a la \emph{estimación
de atenuación}. Su límite reside en que las variables resultantes rara
vez se vinculan de manera simultánea con indicadores de QoS de
aplicación medidos en el mismo experimento.

\subsubsection{Enfoque 2: Medición experimental del desempeño Wi-Fi y de
los protocolos de
aplicación}\label{enfoque-2-mediciuxf3n-experimental-del-desempeuxf1o-wi-fi-y-de-los-protocolos-de-aplicaciuxf3n}

Las mediciones de campo sobre IEEE 802.11n y 802.11ac en oficinas han
demostrado que las tasas teóricas no se alcanzan en la práctica, que el
rendimiento decae con la distancia y que las interferencias de canal
adyacente degradan severamente el enlace (Fiehe y otros, 2010; Dianu y
otros, 2014; Miaoudakis y otros, 2009; van Bloem y otros, 2012). En
paralelo, los estudios experimentales de EDCA han mostrado la influencia
medible de sus parámetros sobre el retardo en plataformas reales
(Parastar y Wang, 2020). Sobre la capa de aplicación, las pruebas de
estrés y los bancos de ensayo con MQTT han establecido el ordenamiento
cualitativo de sus niveles de servicio en términos de retardo, jitter,
pérdida y sobrecarga, tanto en redes de sensores como sobre Wi-Fi con
nodos ESP32 (Hijazi y otros; Akbar y otros, 2017; Nikfarjam Nashtifani,
2026; Chen y Kunz). Este enfoque entrega, así, la \emph{línea base
métrica} de la investigación: sabe cómo se comporta el protocolo
aisladamente. Su límite es que las mediciones se reportan como funciones
de una variable a la vez —distancia, carga o nivel de QoS— y no como un
diseño factorial que permita estimar interacciones, ni acompañadas de un
modelo de predicción.

\subsubsection{Enfoque 3: Predicción y clasificación del estado del
enlace mediante Machine
Learning}\label{enfoque-3-predicciuxf3n-y-clasificaciuxf3n-del-estado-del-enlace-mediante-machine-learning}

La tercera línea concentra los avances más recientes. Se han aplicado
redes neuronales profundas a la evaluación de parámetros QoS de MQTT
(Pawar y otros, 2023), modelos de series temporales y
\emph{transformers} a la predicción de indicadores de red (Hameed,
2023), métodos de ensamble a la clasificación de paquetes y al
enrutamiento asistido por aprendizaje (Çoğay y otros, 2024; Ávalos-Arce
y otros, 2024; Yazici y Güreş, 2024) y aprendizaje profundo a la
anticipación de congestión (Benadji y otros, 2023). Los trabajos
coinciden en tres hallazgos transferibles: los métodos de ensamble
superan consistentemente a los métodos lineales en tareas de predicción
de indicadores; la selección rigurosa de hiperparámetros con validación
cruzada anidada mejora la robustez de las estimaciones; y atributos
asociados a la fortaleza de enlace y a la configuración de transmisión
figuran entre los predictores dominantes. El enfoque de predicción de
calidad de canal con validación experimental confirma además que existe
una brecha entre el desempeño obtenido en condiciones de laboratorio y
el operativo (Becvar y otros, 2024). Su límite es que, en la gran
mayoría de los casos, los predictores provienen de simuladores, de
rastreadores de paquetes o de registros de una sola capa, y no de un
instrumento propio que registre simultáneamente radiofrecuencia y
aplicación en un diseño experimental controlado.

\subsubsection{Enfoque 4: Monitoreo, visualización y gestión de la
calidad de
servicio}\label{enfoque-4-monitoreo-visualizaciuxf3n-y-gestiuxf3n-de-la-calidad-de-servicio}

El cuarto frente comprende los trabajos que trasladan las mediciones a
la operación. Se han propuesto marcos de monitoreo de plataformas IoT
con arquitecturas de streaming y visualización integrada (Calderon y
otros, 2023), marcos de monitoreo de QoS en redes industriales híbridas
orientados a alimentar a los sistemas de gestión (Ehrlich y otros, 2019)
y, en una vertiente de menor complejidad pero alta aplicabilidad,
sistemas de captura de tráfico Wi-Fi con visualización interactiva en
Power BI para el análisis de uso y detección de anomalías en
instalaciones educativas (Praja y otros, 2025). Este enfoque aporta el
mecanismo de \emph{traducción} entre el resultado analítico y la
decisión, pero habitualmente se limita a visualizar datos ya procesados,
sin corresponder a un modelo predictivo validado ni a un criterio de
éxito formalizado.

\subsubsection{Párrafo de vacío}\label{puxe1rrafo-de-vacuxedo}

Del recorrido precedente se desprende que la literatura disponible
ofrece, de manera separada, (i) modelos de propagación en interiores
calibrables pero desconectados de las métricas de calidad de servicio;
(ii) mediciones experimentales del comportamiento Wi-Fi y de los niveles
de servicio MQTT que carecen de un diseño factorial capaz de estimar las
interacciones entre propagación, carga y confiabilidad del protocolo;
(iii) modelos de Machine Learning que predicen indicadores de red sobre
datos de una sola capa o provenientes de simulación, sin instrumentación
propia que registre simultáneamente las variables físicas y las
variables dinámicas; y (iv) tableros de visualización que comunican
descriptores de la red sin sustento en un modelo predictivo ni en un
criterio de éxito formal. En consecuencia, \textbf{existe un vacío de
investigación en la integración experimental de los cuatro frentes: no
se dispone, hasta donde alcanza la revisión efectuada, de un prototipo
IoT híbrido de bajo costo que, mediante una única campaña de mediciones
estructurada en un diseño factorial 3×3×3, registre de forma
sincronizada las variables físicas de radiofrecuencia y las variables
dinámicas de red; correlacione ambas familias mediante modelos de
aprendizaje automático validados; verifique un criterio de éxito
explícito del 85\% con un plan de contingencia predefinido; y comunique
el resultado en tableros de decisión.} La presente investigación se
sitúa exactamente en ese punto de intersección.

\textbf{NOTA:} Si bien los trabajos descritos en los cuatro enfoques
constituyen precedentes legítimos de las técnicas empleadas, la
propuesta se diferencia de ellos de manera sustancial. Frente a los
estudios de propagación (Enfoque 1), no se limita a estimar atenuación,
sino que utiliza la caracterización del canal como fuente de predictores
de una variable de QoS. Frente a las mediciones de protocolo (Enfoque
2), no evalúa un factor aislado, sino que emplea un diseño factorial
completo 3×3×3 que permite estimar efectos principales e interacciones
entre escenario de propagación, carga de tráfico y nivel de servicio
MQTT. Frente a los trabajos de Machine Learning (Enfoque 3), no consume
datos de simuladores ni de una sola capa, sino que construye su propio
instrumento de adquisición sincronizada y somete el modelo a un criterio
de éxito del 85\% regido por un plan de contingencia de cinco pasos.
Frente a los trabajos de visualización (Enfoque 4), no presenta tableros
descriptivos, sino una capa de indicadores alimentada por resultados
estadísticos y predictivos previamente validados. En suma, lo que se
propone hacer no es la repetición de alguno de los enfoques previos,
sino su integración en un único protocolo experimental reproducible.

{[}INSERTAR AQUÍ FLUJOGRAMA DE: Modelo de pila de protocolos y la
interacción cross-layer entre Wi-Fi IEEE 802.11 y MQTT{]}

\subsection{2.4. Marco Conceptual}\label{marco-conceptual}

\textbf{Calidad de Servicio (QoS).} Conjunto de efectos de un servicio
que condicionan la satisfacción del usuario (ITU-T E.800); en este
trabajo, su expresión operacional corresponde al conjunto de métricas de
latencia, jitter, PDR, pérdida y buen uso.

\textbf{Latencia extremo a extremo.} Tiempo transcurrido desde la
emisión de un mensaje por el nodo publicador hasta su recepción efectiva
por el suscriptor, incluyendo acceso al medio, cola en el \emph{broker}
y procesamiento de extremos.

\textbf{Jitter.} Variación de los retardos de entrega consecutivos; se
expresa como desviación típica de la serie de retardos o como diferencia
de retardo de paquetes.

\textbf{Razón de entrega de tramas (PDR).} Porcentaje de mensajes
recibidos correctamente respecto del total de mensajes publicados;
métrica de referencia para el umbral del 85\%.

\textbf{Buen uso (\emph{goodput}).} Carga útil entregada por unidad de
tiempo, excluidas cabeceras, confirmaciones y retransmisiones.

\textbf{Banda ISM.} Bandas de radio Industrial, Científica y Médica, de
uso libre sin licencia, en las que operan Wi-Fi, Bluetooth y Zigbee bajo
la condición de no causar interferencias perjudiciales.

\textbf{Wi-Fi / IEEE 802.11.} Familia de estándares de las capas física
y de enlace para redes inalámbricas locales; en este trabajo, la
variante operativa en 2,4 GHz sobre la que se transmiten las
publicaciones MQTT.

\textbf{DCF y CSMA/CA.} Función de coordinación distribuida y método de
acceso al medio por detección de portadora con evitación de colisiones,
base del mecanismo de acceso de IEEE 802.11.

\textbf{Retroceso exponencial binario.} Procedimiento de diferimiento
aleatorio de la transmisión mediante una ventana de contienda que se
duplica tras cada colisión, hasta un límite máximo.

\textbf{EDCA / WMM.} Mecanismo de acceso mejorado definido por IEEE
802.11e que diferencia cuatro categorías de tráfico con parámetros de
contienda distintos; prioriza, pero no garantiza, calidad de servicio.

\textbf{Nodo no visible (\emph{hidden node}).} Situación topológica en
la que dos estaciones no perciben entre sí pero transmiten hacia el
mismo punto de acceso, provocando colisiones no detectables por el
sondeo de portadora.

\textbf{Multipropagación.} Llegada al receptor de múltiples copias de la
señal por reflexión, difracción y dispersión, origen del desvanecimiento
selectivo.

\textbf{Desvanecimiento por sombra.} Variación lenta de la potencia
recibida debida a obstáculos entre transmisor y receptor; modelada como
variable aleatoria gaussiana en decibelios.

\textbf{Exponente de pérdida de trayecto.} Parámetro n que cuantifica la
rapidez con que la potencia decrece con la distancia en el modelo de
pérdida logarítmica; en interiores su valor depende del material y de la
geometría del sitio.

\textbf{RSSI.} Indicador de intensidad de señal recibida expuesto por el
controlador Wi-Fi; predictor físico principal del prototipo.

\textbf{SNR.} Relación señal-ruido en el receptor; condiciona la tasa de
error alcanzable y, por tanto, la selección de modulación y la necesidad
de reintentos.

\textbf{MQTT.} Protocolo de mensajería ligero basado en
publicación–suscripción, definido por OASIS, empleado para transportar
las mediciones de los nodos hasta el \emph{broker}.

\textbf{Tema (\emph{topic}).} Cadena jerárquica a la que se publican o
de la que se suscriben los mensajes en MQTT; estructura lógica que
organiza los flujos de datos del prototipo.

\textbf{Nivel de QoS MQTT.} Grado de confiabilidad acordado entre
cliente y \emph{broker}: QoS 0 sin confirmación, QoS 1 con confirmación
y QoS 2 con ciclo completo de cuatro manos.

\textbf{\emph{Broker} MQTT.} Servidor que mantiene las sesiones, enruta
los mensajes entre publicadores y suscriptores y gestiona las
confirmaciones.

\textbf{ESP32.} Microcontrolador de bajo costo con interfaz Wi-Fi IEEE
802.11 y Bluetooth integrado; plataforma del nodo instrumentador del
prototipo.

\textbf{Prototipo IoT híbrido.} Conjunto hardware–software que integra
nodos ESP32, infraestructura Wi-Fi, \emph{broker} MQTT, base de datos y
capa de análisis para la adquisición conjunta de variables físicas y de
red.

\textbf{Aprendizaje supervisado.} Paradigma de Machine Learning en el
que un modelo se ajusta a un conjunto de datos etiquetado para predecir
una variable objetivo.

\textbf{Clasificador / regresor.} Modelo que predice una categoría
discreta / una magnitud continua.

\textbf{Bosque aleatorio (\emph{random forest}).} Método de ensamble que
combina múltiples árboles de decisión entrenados sobre muestras y
subconjuntos aleatorios de atributos.

\textbf{Validación cruzada en \emph{k} pliegues.} Procedimiento de
evaluación en el que el conjunto de datos se divide en \emph{k} partes y
el modelo se entrena y evalúa en rotación, estimando la capacidad de
generalización.

\textbf{Exactitud, sensibilidad, precisión y \emph{F}1.} Métricas de
desempeño de un clasificador; \emph{F}1 es la media armónica entre
precisión y sensibilidad, recomendable ante clases desequilibradas.

\textbf{Diseño factorial experimental.} Plan de experimentación en el
que se manipulan simultáneamente varios factores con sus respectivos
niveles, permitiendo estimar efectos principales e interacciones.

\textbf{Efecto principal e interacción.} Influencia promedio de un
factor sobre la respuesta; influencia conjunta de dos o más factores que
no resulta de la suma de sus efectos individuales.

\textbf{Umbral del 85\%.} Criterio de éxito metodológico
operacionalizado: el porcentaje mínimo de cumplimiento que debe alcanzar
la métrica de referencia (PDR) y la exactitud del modelo predictivo para
considerar verificada la hipótesis de trabajo.

\textbf{Plan de contingencia de cinco pasos.} Procedimiento predefinido
de respuesta ante resultados que no alcanzan el umbral del 85\%, con
pasos de verificación instrumental, re-aislamiento de causas, ampliación
de evidencia, reajuste de parámetros y documentación formal de la
desviación.

\textbf{Power BI.} Plataforma de inteligencia de negocio para la
modelación de datos, el cálculo de medidas y la construcción de tableros
interactivos con indicadores de desempeño.

\textbf{KPI y tablero.} Indicador clave de desempeño y superficie visual
que lo comunica con tendencias, metas y estados de cumplimiento.

\subsection{2.5. Normatividad y
Legislación}\label{normatividad-y-legislaciuxf3n}

La normatividad aplicable se limita a los documentos estrictamente
pertinentes con el objeto de estudio:

\textbf{Estándares IEEE 802.11.} El estándar IEEE 802.11 y sus
modificaciones (802.11n, 802.11e/EDCA, 802.11ac, 802.11ax) definen las
reglas de convivencia del enlace utilizado por el prototipo:
modulaciones, canales, mecanismos de acceso al medio y parámetros de
contienda. Su cumplimiento es condición de que las mediciones sean
comparables con las de terceros.

\textbf{Estándar OASIS MQTT.} La especificación MQTT Version 5.0 y su
predecesora 3.1.1, publicadas por el consorcio OASIS, definen la
semántica de los niveles de servicio, los flujos de confirmación y los
parámetros de sesión cuyo efecto se mide experimentalmente.

\textbf{Recomendaciones ITU-T.} La Recomendación E.800 aporta la
definición normativa de calidad de servicio; la Recomendación G.114 fija
las referencias de tolerancia al retardo de ida; la serie Y.1540
establece clases de desempeño de transferencia de paquetes IP que se
toman como marco interpretativo de los resultados.

\textbf{Recomendaciones ITU-R.} La Recomendación ITU-R P.1238
proporciona los modelos y parámetros de propagación en interiores
empleados como línea base de comparación del modelo calibrado en el
escenario de medición.

\textbf{Reglamento de Radiocomunicaciones de la Unión Internacional de
Telecomunicaciones.} El Reglamento establece el régimen de las bandas
ISM, de uso libre y sin licencia, bajo la condición de no causar
interferencias perjudiciales a otros servicios; esta condición delimita
el marco de operación del prototipo y de los equipos vecinos con los que
coexiste.

\textbf{Marco de protección de datos personales.} La normativa nacional
de protección de datos personales resulta pertinente en tanto la
investigación opera con registros de tráfico y de estado de enlace. El
protocolo experimental contempla explícitamente la recolección exclusiva
de métricas de red y de radiofrecuencia, sin captura de contenido de las
comunicaciones ni de datos personales de los usuarios, condición que se
verifica durante la etapa de instrumentación y se declara en el informe
final.

\textbf{Aspectos de seguridad.} Si bien la seguridad del transporte no
constituye una variable de la investigación, la utilización de canales
cifrados y de credenciales de acceso al \emph{broker} se rige por las
prácticas recomendadas en las propias especificaciones de OASIS MQTT y
por la configuración segura del punto de acceso inalámbrico, con el
propósito de que el tráfico experimental no introduzca sesgos ni
vulnerabilidades en la red de la institución anfitriona.

\section{CAPÍTULO III: ELABORACIÓN DE LA
PROPUESTA}\label{capuxedtulo-iii-elaboraciuxf3n-de-la-propuesta}

El presente capítulo describe la propuesta en su integridad: las
generalidades del diseño de investigación (3.1), el esquema del
prototipo (3.2), la definición operacional de las variables (3.3), el
diseño factorial experimental 3×3×3 que estructura la campaña de
mediciones (3.4), el criterio de éxito del 85\% con su plan de
contingencia explícito de cinco pasos (3.5), el alcance de la
intervención (3.6) y el procedimiento de validación que garantiza el
rigor científico de los resultados (3.7).

\subsection{3.1. Generalidades}\label{generalidades}

\subsubsection{3.1.1. Tipo, nivel, enfoque y método de
investigación}\label{tipo-nivel-enfoque-y-muxe9todo-de-investigaciuxf3n}

\textbf{Tipo.} La investigación es \emph{aplicada}, porque su finalidad
inmediata es producir un artefacto —el prototipo IoT híbrido— y un
conjunto de reglas de uso —el protocolo de evaluación del QoS— que
resuelven una problemática concreta de ingeniería. Es también de
carácter \emph{experimental}, porque manipula deliberadamente factores
de un sistema y observa la respuesta bajo condiciones controladas, lo
que habilita la atribución causal de los cambios observados.

\textbf{Nivel.} El nivel es \emph{explicativo-correlacional}: la primera
fase explica la influencia de los factores experimentales sobre las
métricas de calidad de servicio mediante análisis de varianza; la
segunda establece la relación estadística y predictiva entre las
variables físicas de radiofrecuencia y las variables dinámicas de red
mediante modelos de aprendizaje automático. El nivel explicativo se
justifica por el carácter manipulado de los factores; el correlacional,
por el carácter observacional de los índices de radiofrecuencia, que no
se manipulan sino que se miden.

\textbf{Enfoque.} Cuantitativo, con uso secundario de información de
carácter cualitativo únicamente en la documentación de incidentes
operativos (observaciones del laboratorio, notas de campo sobre
condiciones ambientales) que se emplean como contexto interpretativo y
no como evidencia estadística.

\textbf{Método.} Método experimental con dos momentos complementarios:
(i) \emph{medición y comparación} bajo un diseño factorial completo, que
responde a la PE2; y (ii) \emph{modelado predictivo supervisado}, que
responde a la PE3. Ambos momentos comparten la misma base de datos,
generada por un único instrumento de adquisición, condición que
garantiza la coherencia entre el análisis de efectos y el análisis de
correlación.

\subsubsection{3.1.2. Lugar de ejecución}\label{lugar-de-ejecuciuxf3n}

La investigación se ejecuta en el interior de una casa habitación de uso
residencial, acondicionada como escenario experimental controlado. La
elección de un entorno residencial obedece a tres razones de validez
ecológica: (i) constituye el escenario natural del despliegue IoT
doméstico, que es el horizonte de aplicación declarado en la PE4, de
modo que las condiciones de propagación observadas son representativas
del contexto real de uso; (ii) ofrece una geometría de recintos cerrados
con tabiquería de concreto, mobiliario y distribución de puertas y
ventanas equivalente a la de cualquier vivienda urbana estándar, lo que
permite caracterizar atenuaciones típicas de muros de carga y tabiques
no estructurales; y (iii) al estar bajo control exclusivo del
experimentador, permite fijar la ubicación del punto de acceso, de los
nodos y del \emph{broker}, y mantener constantes las condiciones de
ocupación del espacio durante los ensayos.

El escenario se encuentra delimitado por un plano o croquis de planta de
la vivienda, levantado a escala aproximada y con las distancias entre
los puntos de emplazamiento medidas en metros. Este plano cumple una
función metodológica triple: (i) sirve como instrumento de trazabilidad,
pues fija de manera permanente y reproducible las coordenadas de
emplazamiento del punto de acceso, de los tres nodos emisores (N1, N2 y
N3), del nodo suscriptor y de la estación de adquisición; (ii)
constituye el soporte sobre el cual se codifican los niveles del factor
A —escenario de propagación— conforme a la trayectoria radioeléctrica
que media entre emisor y receptor; y (iii) permite a un tercero
reconstruir la geometría del experimento sin necesidad de acudir al
escenario físico, condición exigida por la validación experimental
descrita en el apartado 3.7. El croquis se incorpora al cuaderno de
campo y a los anexos del documento final, junto con fotografías del
emplazamiento correspondientes a cada nivel de propagación.

Sobre la base de este levantamiento geométrico, los tres niveles del
factor A se operacionalizan en el entorno residencial de la siguiente
manera:

\begin{itemize}
\item
  \textbf{A1 — Visibilidad directa.} El nodo emisor y el nodo suscriptor
  se disponen dentro del mismo recinto o en recintos contiguos
  comunicados por puerta abierta, con línea visual directa entre ambos y
  sin muro de concreto interpuesto. La distancia entre emisor y receptor
  se fija en un rango corto, de orden métrico, representativo de la
  proximidad típica entre dispositivos IoT domésticos —por ejemplo, un
  sensor y su pasarela en la misma habitación o en habitaciones
  contiguas comunicadas por espacio abierto—.
\item
  \textbf{A2 — Muro de concreto.} El nodo suscriptor se traslada a un
  recinto distinto del emisor, de modo que el trayecto radioeléctrico
  atraviesa exactamente un muro de concreto —de carga o tabique
  equivalente— sin visibilidad directa entre ambos puntos. La distancia
  geométrica se mantiene en un orden de magnitud comparable al del nivel
  A1, de manera que la diferencia observada entre A1 y A2 se atribuya
  predominantemente a la atenuación del muro y no a la pérdida por
  distancia, en concordancia con el modelo de pérdida de potencia
  multi-muro desarrollado en el marco teórico.
\item
  \textbf{A3 — Doble muro o ruido ISM elevado.} El nodo suscriptor se
  ubica en un recinto cuya comunicación con el emisor requiere atravesar
  dos muros de concreto —por ejemplo, entre extremos opuestos de la
  vivienda o entre el área principal y un espacio semiabierto contiguo—
  o, alternativamente, el ensayo se ejecuta en un periodo en el cual la
  banda ISM de 2,4 GHz presenta un ruido de fondo elevado por
  interferencia de redes vecinas o de aparatos electrodomésticos en
  operación. En cualquiera de las dos variantes, el nivel A3 representa
  la condición de propagación severa; la variante empleada en cada
  ensayo se registra explícitamente en el cuaderno de campo para efectos
  de trazabilidad y de interpretación de los resultados.
\end{itemize}

El punto de acceso —router Wi-Fi doméstico— y la estación de
adquisición, sobre la cual se instala el \emph{broker} Mosquitto, se
emplazan en posiciones fijas y de antemano definidas en el croquis, que
permanecen inalterables durante toda la campaña. El \emph{broker} MQTT y
la base de datos se operan en una red local (LAN) del propio escenario
experimental, de modo que el trayecto de cada mensaje quede íntegramente
bajo control del experimentador y la influencia de Internet sobre la
latencia quede excluida de las mediciones. Para ello, el router
permanece operativo en modo aislado —sin enlace con Internet o con la
WAN desconectada— durante toda la campaña principal, de forma que el
tráfico externo no introduzca una fuente de variación no controlada
sobre el medio compartido de la banda ISM; únicamente los ensayos
expresamente destinados a evaluar el escenario remoto, ejecutados en una
etapa complementaria y fuera de la matriz factorial principal, se
realizan con el enlace a Internet habilitado.

Las condiciones ambientales del escenario —ocupación humana, estado de
puertas y ventanas, aparatos electrodomésticos encendidos o apagados y
presencia o ausencia de personas en los recintos— se fijan de antemano y
se mantienen constantes dentro de cada ensayo, registrándose en el
cuaderno de campo toda desviación. Este control ambiental es
indispensable en un entorno residencial, donde la presencia humana y el
funcionamiento de electrodomésticos constituyen fuentes de atenuación y
de interferencia cuya variación no controlada degradaría la validez
interna del experimento.

\subsubsection{3.1.3. Material y equipos}\label{material-y-equipos}

{\def\LTcaptype{none} % do not increment counter
\begin{longtable}[]{@{}
  >{\raggedright\arraybackslash}p{(\linewidth - 4\tabcolsep) * \real{0.3333}}
  >{\raggedright\arraybackslash}p{(\linewidth - 4\tabcolsep) * \real{0.3333}}
  >{\raggedright\arraybackslash}p{(\linewidth - 4\tabcolsep) * \real{0.3333}}@{}}
\toprule\noalign{}
\begin{minipage}[b]{\linewidth}\raggedright
Componente
\end{minipage} & \begin{minipage}[b]{\linewidth}\raggedright
Especificación funcional
\end{minipage} & \begin{minipage}[b]{\linewidth}\raggedright
Rol en el experimento
\end{minipage} \\
\midrule\noalign{}
\endhead
\bottomrule\noalign{}
\endlastfoot
Nodos emisores (3 unidades) & Placa ESP32 NodeMCU DevKit V1 con interfaz
Wi-Fi IEEE 802.11 b/g/n, firmware propio con biblioteca MQTT y lectura
de RSSI/SNR del controlador & Generan la carga de tráfico y reportan los
índices de radiofrecuencia en cada publicación \\
Nodo suscriptor (1 unidad) & Placa ESP32 NodeMCU DevKit V1 con conexión
Wi-Fi y marca de tiempo local & Registra la recepción, calcula la
latencia extremo a extremo y verifica la unicidad de los mensajes \\
Punto de acceso inalámbrico & Router Wi-Fi doméstico operativo en la
banda ISM de 2,4 GHz, en un canal fijo no solapado (1, 6 u 11) & Define
el medio de acceso y las condiciones de contienda \\
\emph{Broker} MQTT & \emph{Broker} de código abierto (por ejemplo,
Mosquitto) instalado en la estación de adquisición, con registro
(\emph{logging}) de sesión y de mensajes & Enruta y confirma los
mensajes según el nivel de QoS ensayado \\
Estación de adquisición & Computadora con sistema operativo, servidor de
base de datos y sincronización NTP & Almacena las series temporales
sincronizadas de ambas familias de variables \\
Estación de análisis & Computadora con entorno de programación Python,
bibliotecas de análisis de datos, Scikit-learn y Power BI & Ejecuta el
procesamiento, el modelado y la construcción de tableros \\
Instrumentación complementaria & Medidor de campo o aplicación de
medición Wi-Fi para verificación independiente del RSSI; plancha de
planta del escenario con distancias marcadas & Verificación cruzada del
instrumento y levantamiento geométrico del escenario \\
\end{longtable}
}

La selección del hardware obedece al criterio de costo por nodo
compatible con despliegues masivos, al que se alude en la PE4; la lista
de materiales y su presupuesto corresponden a los anexos del documento
final.

{[}INSERTAR AQUÍ TABLA: Especificación técnica y costos de los
componentes del prototipo físico ESP32{]}

{[}INSERTAR AQUÍ IMAGEN TÉCNICA: Topología de nodos ESP32 — nodos
emisores N1–N3, suscriptor y punto de acceso en el escenario
experimental{]}

\subsubsection{3.1.4. Procedimientos}\label{procedimientos}

La ejecución se organiza en cinco fases secuenciales, con entregables
verificables en cada una de ellas:

\begin{enumerate}
\def\labelenumi{\arabic{enumi}.}
\tightlist
\item
  \textbf{Diseño y construcción del prototipo.} Definición de la
  arquitectura de mensajes, codificación de los campos de la publicación
  (identificador de nodo, marca de tiempo de emisión, contador de
  secuencia, RSSI, SNR, nivel de ruido, escenario y nivel de QoS),
  integración del \emph{broker} y de la base de datos, y pruebas de
  integración en mesa.
\item
  \textbf{Calibración del instrumento.} Verificación de la
  sincronización de relojes mediante NTP, medición de referencia en un
  trayecto conocido para el ajuste del modelo de pérdida de distancia, y
  validación cruzada del RSSI reportado por el ESP32 con un instrumento
  independiente. Criterio de salida: error de sincronización y de
  medición dentro de los tolerancias declarados en el protocolo.
\item
  \textbf{Campaña de mediciones.} Ejecución de la matriz factorial
  completa con el número de réplicas predefinido, en orden aleatorizado
  de ensayos, registrando todas las condiciones y todo incidente en el
  cuaderno de campo.
\item
  \textbf{Procesamiento y modelado.} Tratamiento estadístico de los
  datos, verificación de supuestos, análisis de varianza, entrenamiento
  y validación de los modelos de Machine Learning con Scikit-learn.
\item
  \textbf{Comunicación de resultados y validación.} Construcción de los
  tableros de Power BI, contrastación con el criterio de éxito,
  aplicación del plan de contingencia cuando corresponda y redacción de
  conclusiones.
\end{enumerate}

{[}INSERTAR AQUÍ FLUJOGRAMA DE: Las cinco fases secuenciales de
ejecución de la investigación con sus entregables verificables{]}

\subsubsection{3.1.5. Protocolo operativo de la campaña de
mediciones}\label{protocolo-operativo-de-la-campauxf1a-de-mediciones}

El protocolo operativo fija, de manera previa e inmutable durante la
campaña, las condiciones bajo las cuales se ejecuta cada ensayo del
diseño factorial 3×3×3. Su objetivo es garantizar que las 27 celdas del
diseño se midan bajo condiciones homogéneas y comparables, y que la
variabilidad observada en las respuestas sea atribuible a los factores
ensayados y no a diferencias de procedimiento.

\textbf{Equipamiento del escenario.} El equipamiento se compone de: (i)
cuatro placas ESP32 NodeMCU DevKit V1 —tres configuradas como nodos
emisores N1, N2 y N3 y una como nodo suscriptor—, cargadas con un
firmware común que incluye la pila TCP/IP, el cliente MQTT y la lectura
de los índices de radiofrecuencia (RSSI, SNR y nivel de ruido de fondo)
reportados por el controlador en cada publicación; (ii) un router Wi-Fi
doméstico operativo en la banda ISM de 2,4 GHz, configurado en un canal
fijo no solapado —1, 6 u 11, según el verificado mediante análisis de
espectro previo—, que actúa como punto de acceso y medio de contienda;
(iii) una estación de adquisición —computadora de escritorio o portátil—
sobre la cual se instalan el \emph{broker} Mosquitto y la base de datos
que recibe, enruta y almacena las publicaciones; y (iv) la
instrumentación complementaria descrita en el apartado 3.1.3. Todos los
equipos se emplazan conforme a las coordenadas fijadas en el croquis de
la vivienda y no se desplazan durante la ejecución de una celda.

\textbf{Condiciones de ensayo.} Cada ensayo se ejecuta bajo las
siguientes condiciones fijas, verificadas antes de su inicio y
registradas en el cuaderno de campo: sincronización NTP activa y
verificada en los cuatro nodos y en la estación de adquisición, con el
desfase observado registrado; tamaño de carga útil constante e idéntico
en todos los nodos; firmware idéntico y de la misma versión compilada en
las cuatro placas; estado del entorno residencial conforme a las
condiciones ambientales fijadas en el apartado 3.1.2; y nivel de
ocupación humana del escenario nulo o estrictamente constante. Se
verifica además que el canal Wi-Fi seleccionado no haya sido alterado
por reinicios o reconexiones del router desde la calibración del
instrumento.

\textbf{Duración y volumen de cada ensayo.} Cada ensayo tiene una
duración fija de diez (10) minutos, contados desde el inicio de la
publicación programada del primer nodo emisor hasta la recepción del
último mensaje del nodo suscriptor o el cierre de la sesión MQTT en el
\emph{broker}, lo que ocurra primero. Esta duración se fija con base en
dos consideraciones: (i) es suficientemente larga para acumular un
volumen de mensajes representativo incluso en el nivel de carga más
esparsa, de modo que las métricas de proporción (PDR y pérdida de
paquetes) se calculen sobre una base muestral estadísticamente
admisible; y (ii) es suficientemente corta para que las condiciones
ambientales y de interferencia se mantengan razonablemente estables
dentro del ensayo, reduciendo la influencia de tendencias temporales
intra-ensayo. La frecuencia de publicación de cada nodo emisor la
determina el nivel del factor B vigente en la celda —B1 oferta esparsa,
B2 oferta moderada, B3 oferta intensa—, definido operacionalmente por el
intervalo entre publicaciones consecutivas del mismo nodo, cuyos valores
exactos en milisegundos se fijan antes de la campaña y se reportan en el
capítulo de resultados.

\textbf{Réplicas y orden aleatorizado.} Cada una de las 27 celdas del
diseño factorial se ejecuta con tres (3) réplicas independientes, lo que
arroja un total de 81 ensayos para la campaña completa. La repetición en
tres réplicas se justifica por la necesidad de estimar la varianza
intra-celda —requisito del análisis de varianza previsto en el apartado
3.4.4— y de verificar la reproducibilidad de las métricas dentro de cada
condición experimental. El orden de ejecución de los 81 ensayos se
determina mediante una secuencia aleatoria generada antes del inicio de
la campaña, con el fin de neutralizar el efecto de posibles tendencias
temporales —variación de la interferencia de redes vecinas, cambios de
temperatura o de ocupación del entorno residencial a lo largo del día—
sobre las respuestas medidas. La secuencia aleatoria se registra en el
cuaderno de campo y se conserva como evidencia de la aleatorización
efectiva. Se prohíbe la reordenación de ensayos por conveniencia
operativa una vez iniciada la campaña; si un ensayo debe suspenderse por
causa de fuerza mayor, se reprograma en la primera ventana disponible y
la suspensión se documenta.

\textbf{Registro y trazabilidad de cada ensayo.} Cada ensayo genera tres
registros triangulados —en el nodo emisor (marca de tiempo de emisión,
RSSI, SNR, nivel de ruido y contador de secuencia), en el nodo
suscriptor (marca de tiempo de recepción, contador de secuencia y
verificación de unicidad) y en el \emph{broker} Mosquitto (sesión,
suscripción, mensaje recibido y confirmación de entrega conforme al
nivel de QoS ensayado)—, los cuales se consolidan en la base de datos
bajo una etiqueta única compuesta por el identificador de celda
(AᵢBⱼCₖ), el número de réplica (1, 2 o 3) y el identificador de
secuencia del mensaje. A cada ensayo se le asocian, además de los tres
registros triangulados, la marca de tiempo de inicio y de fin, la
versión del firmware cargado en las placas, la configuración vigente del
router (canal, potencia de transmisión, identificador SSID y modo de
seguridad), la versión del \emph{broker} Mosquitto, la condición
ambiental del escenario (estado de puertas, ocupación y aparatos
electrodomésticos) y cualquier incidente operativo registrado por el
experimentador. Esta información se asienta en el cuaderno de campo,
instrumento primario de trazabilidad de la campaña cuyo formato se
describe en los anexos del documento final.

\textbf{Condiciones de parada y reanudación.} Un ensayo se interrumpe y
se anula —sin contabilizarse como réplica válida— cuando se produce
cualquiera de los siguientes eventos: (i) desconexión o reinicio no
programado del router Wi-Fi; (ii) pérdida de sincronización NTP
verificable en cualquiera de los cuatro nodos; (iii) interrupción del
servicio del \emph{broker} Mosquitto o de la base de datos en la
estación de adquisición; (iv) alteración no programada del entorno
residencial —entrada de ocupantes, apertura o cierre de puertas,
encendido o apagado de aparatos electrodomésticos relevantes— que
comprometa la constancia de las condiciones ambientales fijadas. En caso
de interrupción, se documenta la causa en el cuaderno de campo, se
restablecen las condiciones de ensayo y se reejecuta el ensayo íntegro
desde su inicio, en la primera ventana disponible conforme a la
secuencia aleatorizada. Ningún ensayo interrumpido se contabiliza
parcialmente: la integridad de cada réplica es condición indispensable
de su validez.

\subsection{3.2. Esquema de la Propuesta}\label{esquema-de-la-propuesta}

La propuesta se compone de cuatro capas funcionales que se interconectan
de la siguiente manera:

\begin{verbatim}
 CAPA 1 — CAPTACIÓN Y EMISIÓN (nodos instrumentados)
 ┌────────────────────┐   ┌────────────────────┐   ┌────────────────────┐
 │  Nodo ESP32 (N1)   │   │  Nodo ESP32 (N2)   │   │  Nodo ESP32 (N3)   │
 │  + marca de tiempo │   │  + marca de tiempo │   │  + marca de tiempo │
 │  + RSSI / SNR      │   │  + RSSI / SNR      │   │  + RSSI / SNR      │
 └─────────┬──────────┘   └─────────┬──────────┘   └─────────┬──────────┘
           └───────────────┬────────┴────────────┬───────────┘
                           ▼                     ▼
 CAPA 2 — MEDIO INALÁBRICO     ┌──────────────────────────────┐
 (Wi-Fi IEEE 802.11, 2,4 GHz)  │ Punto de acceso (canal fijo  │
 escenarios A1 / A2 / A3       │ no solapado, CSMA/CA, DCF)   │
                               └──────────────┬───────────────┘
                                              ▼
 CAPA 3 — TRANSPORTE Y ALMACENAMIENTO
                    ┌────────────────────────────────────────┐
                    │  Broker MQTT (QoS 0 / 1 / 2 = C1/C2/C3) │
                    └────────────────────┬───────────────────┘
                                         ▼
                    ┌────────────────────────────────────────┐
                    │  Estación de adquisición + base de      │
                    │  datos (series sincronizadas por NTP)   │
                    └────────────────────┬───────────────────┘
                                         ▼
 CAPA 4 — ANÁLISIS Y COMUNICACIÓN
          ┌──────────────────────────┬───────────────────────────┐
          │ Scikit-learn             │ Power BI                  │
          │ ANOVA 3 factores,        │ tableros de KPI,          │
          │ modelos de clasificación │ estados del umbral 85 %,  │
          │ y regresión, importancia │ tendencias por escenario  │
          │ de predictores           │ y alertas de degradación  │
          └──────────────────────────┴───────────────────────────┘
\end{verbatim}

La lectura del esquema es la siguiente: los nodos de la capa 1 emiten
publicaciones MQTT que contienen, en una misma unidad de registro, la
variable temporal de la aplicación y los índices de radiofrecuencia; la
capa 2 condiciona esa transmisión conforme al escenario de propagación
vigente y a las reglas de acceso al medio; la capa 3 materializa el
nivel de servicio ensayado y almacena el resultado; y la capa 4
transforma el conjunto de registros en evidencia estadística, en modelo
predictivo y en indicadores comunicables. La originalidad de la
propuesta reside en la capa 1: el instrumento de medición es \emph{el
propio nodo de la red}, de modo que la obtención de las variables
físicas no requiere de un analizador de espectro ni de una sonda
comercial.

{[}INSERTAR AQUÍ IMAGEN TÉCNICA: Diagrama de arquitectura del prototipo
con las cuatro capas funcionales — captación, medio inalámbrico,
transporte/almacenamiento y análisis{]}

\subsection{3.3. Variables de la
Investigación}\label{variables-de-la-investigaciuxf3n}

\textbf{Variables independientes (factores experimentales).}

\begin{itemize}
\tightlist
\item
  \textbf{A — Escenario de propagación.} Tres niveles operacionalizados
  en el entorno residencial: A1, visibilidad directa entre emisor y
  receptor dentro del mismo recinto o de recintos contiguos abiertos,
  sin muro interpuesto; A2, sin visibilidad directa con exactamente un
  muro de concreto atravesado por el trayecto radioeléctrico; A3, sin
  visibilidad directa con doble muro de concreto atravesado o con ruido
  de fondo ISM elevado.
\item
  \textbf{B — Carga de ofrecimiento.} Tres niveles definidos por el
  intervalo de publicación: B1, oferta esparsa; B2, oferta moderada; B3,
  oferta intensa.
\item
  \textbf{C — Nivel de servicio MQTT.} Tres niveles: C1 (QoS 0), C2 (QoS
  1) y C3 (QoS 2).
\end{itemize}

La combinación de los tres factores con tres niveles produce un
\textbf{diseño factorial completo 3×3×3 con 27 combinaciones de
tratamiento}.

\textbf{Variables dependientes (respuestas).}

\begin{enumerate}
\def\labelenumi{\arabic{enumi}.}
\tightlist
\item
  PDR, en porcentaje — variable de referencia para el umbral del 85\%.
\item
  Latencia extremo a extremo, en milisegundos (media y percentil 95).
\item
  Jitter, en milisegundos (desviación típica de los retardos).
\item
  Pérdida de paquetes, en porcentaje.
\item
  Buen uso, en kilobits por segundo.
\end{enumerate}

\textbf{Variables medidas o covariables físicas} (no manipuladas,
registradas en cada publicación): RSSI, SNR, nivel de ruido de fondo,
distancia entre emisor y receptor, número y tipo de obstáculos del
trayecto.

\textbf{Variables intervinientes o de control:} densidad de nodos fija
(tres emisores y un suscriptor), canal Wi-Fi fijo no solapado, tamaño de
carga útil constante, firmware idéntico en todos los nodos, ubicación
fija del punto de acceso y del \emph{broker}, sincronización NTP activa,
y horario de ejecución registrado. El control de estas variables es lo
que permite atribuir las diferencias observadas a los tres factores
ensayados.

{[}INSERTAR AQUÍ TABLA: Definición operacional de variables —
independientes (factores A/B/C), dependientes (PDR, latencia, jitter,
pérdida, buen uso), covariables físicas y de control{]}

\subsection{3.4. Diseño Experimental Factorial
3×3×3}\label{diseuxf1o-experimental-factorial-333}

\subsubsection{3.4.1. Justificación del
diseño}\label{justificaciuxf3n-del-diseuxf1o}

El problema descrito en la realidad problemática es de naturaleza
\emph{interactiva}: se sospecha que la degradación del QoS no es
aditiva, es decir, que el efecto de una carga elevada sobre la latencia
no es el mismo en visibilidad directa que tras doble muro, y que el
nivel de servicio MQTT modula esa diferencia. Un diseño de un solo
factor a la vez (\emph{one-factor-at-a-time}) es incapaz de detectar
interacciones y exige un número mayor de ensayos para una cobertura
menor. El diseño factorial completo estima simultáneamente los tres
efectos principales y las tres interacciones de dos factores más la de
tres factores, con eficiencia muestral apreciable, y ofrece además una
estructura natural para el análisis de varianza.

\subsubsection{3.4.2. Matriz experimental}\label{matriz-experimental}

{\def\LTcaptype{none} % do not increment counter
\begin{longtable}[]{@{}
  >{\raggedright\arraybackslash}p{(\linewidth - 4\tabcolsep) * \real{0.3333}}
  >{\raggedright\arraybackslash}p{(\linewidth - 4\tabcolsep) * \real{0.3333}}
  >{\raggedright\arraybackslash}p{(\linewidth - 4\tabcolsep) * \real{0.3333}}@{}}
\toprule\noalign{}
\begin{minipage}[b]{\linewidth}\raggedright
Factor
\end{minipage} & \begin{minipage}[b]{\linewidth}\raggedright
Niveles
\end{minipage} & \begin{minipage}[b]{\linewidth}\raggedright
Descripción
\end{minipage} \\
\midrule\noalign{}
\endhead
\bottomrule\noalign{}
\endlastfoot
A — Escenario de propagación & A1 / A2 / A3 & Visibilidad directa / un
muro de concreto / doble muro o ruido ISM elevado \\
B — Carga de ofrecimiento & B1 / B2 / B3 & Intervalo de publicación
esparsa / moderado / intenso \\
C — Nivel de servicio MQTT & C1 / C2 / C3 & QoS 0 / QoS 1 / QoS 2 \\
\end{longtable}
}

La matriz resultante contiene las 27 celdas AiBjCk. Cada celda se
ejecuta con tres réplicas definidas de forma previa, ordenadas
aleatoriamente en el tiempo para neutralizar el efecto de posibles
tendencias temporales (variación de la interferencia de red vecina a lo
largo del día). Cada réplica corresponde a un ensayo de duración fija de
diez minutos con un número mínimo preestablecido de mensajes publicados
—según lo fijado en el protocolo operativo del apartado 3.1.5—, de modo
que las métricas de proporción (PDR, pérdida) se calculen sobre una base
muestral suficiente.

El volumen total de observaciones resulta de 27 celdas multiplicado por
el número de réplicas; a cada mensaje individual se le asocian, además
de su resultado de entrega, su marca de tiempo, los índices de
radiofrecuencia reportados por el nodo y la etiqueta de celda, lo que
conforma un conjunto de datos tabular apto tanto para el análisis de
varianza como para el aprendizaje automático.

{[}INSERTAR AQUÍ TABLA: Matriz completa del diseño factorial 3×3×3 — 27
celdas A\_iB\_jC\_k con niveles descriptos{]}

{[}INSERTAR AQUÍ FLUJOGRAMA DE: Flujo del diseño factorial 3×3×3 —
factores, niveles, réplicas y orden aleatorizado de ensayos{]}

\subsubsection{3.4.3. Instrumentación de la
medición}\label{instrumentaciuxf3n-de-la-mediciuxf3n}

\begin{itemize}
\tightlist
\item
  \textbf{Sincronización.} Todos los nodos sincronizan su reloj por NTP
  antes de cada ensayo; la latencia de un solo sentido se calcula como
  diferencia de marcas de tiempo entre publicación y recepción. Como
  control de calidad, se registra el desfase de sincronización observado
  en cada sesión.
\item
  \textbf{Registro doble.} Cada publicación transporta la información de
  radiofrecuencia; el suscriptor registra la recepción; el \emph{broker}
  conserva el registro de entrega y de confirmación. La triangulación de
  los tres registros permite distinguir un mensaje perdido en el enlace
  inalámbrico de uno perdido en la cola del \emph{broker}.
\item
  \textbf{Etiquetado.} Cada registro lleva los campos de escenario,
  carga, nivel de QoS, réplica y número de secuencia, lo que permite
  reconstruir íntegramente la matriz experimental desde la base de
  datos.
\item
  \textbf{Trazabilidad.} Se conservan la versión del firmware, la
  configuración del punto de acceso, la versión del \emph{broker}, la
  plancha del escenario con las distancias medidas y el cuaderno de
  campo.
\end{itemize}

\subsubsection{3.4.4. Análisis estadístico
previsto}\label{anuxe1lisis-estaduxedstico-previsto}

El análisis de los datos del diseño se plantea mediante un
\textbf{análisis de varianza de tres vías} con modelo de efectos fijos:

Y = μ + Ai + Bj + Ck + (AB)ij + (AC)ik + (BC)jk + (ABC)ijk + ε

donde Y es la respuesta considerada (PDR, latencia, jitter, pérdida o
buen uso) y ε el error experimental. El procedimiento contempla:

\begin{enumerate}
\def\labelenumi{\arabic{enumi}.}
\tightlist
\item
  Verificación de los supuestos mediante la prueba de Shapiro-Wilk sobre
  los residuos y la prueba de Levene sobre la homogeneidad de varianzas;
  cuando las respuestas de proporción violen los supuestos, se aplicará
  una transformación logarítmica o arcseno sobre los porcentajes, y, en
  su defecto, se recurrirá a pruebas no paramétricas alternativas.
\item
  Estimación de los efectos principales y de las interacciones con un
  nivel de significancia fijado por adelantado en α = 0,05.
\item
  Contrastes de comparación múltiple por el método de Tukey para las
  diferencias entre niveles de un mismo factor, con intervalos de
  confianza reportados.
\item
  Estimación de tamaños de efecto y de la varianza explicada por cada
  término del modelo, de modo que la significancia estadística no se
  interprete desligada de la relevancia práctica.
\item
  Verificación del cumplimiento o incumplimiento del umbral del 85\%
  celda por celda, con el intervalo de confianza correspondiente.
\end{enumerate}

{[}INSERTAR AQUÍ FLUJOGRAMA DE: Procedimiento del ANOVA de tres vías —
verificación de supuestos, efectos e interacciones, contrastes de Tukey
y verificación del umbral{]}

\subsubsection{3.4.5. Proceso de aprendizaje automático sobre la misma
base de
datos}\label{proceso-de-aprendizaje-automuxe1tico-sobre-la-misma-base-de-datos}

Sobre el mismo conjunto de registros, y respondiendo a la PE3, se
ejecutará el siguiente flujo: separación de predictores (índices de
radiofrecuencia, distancia, escenario, carga, nivel de QoS) y de
variable objetivo (estado de cumplimiento del umbral y, en la vertiente
de regresión, los valores de latencia y jitter); partición estratificada
entre entrenamiento y prueba; preprocesamiento con imputación y escalado
cuando el algoritmo lo requiera; ajuste de una línea base de regresión
regularizada y de modelos de árbol y ensamble; búsqueda de
hiperparámetros mediante validación cruzada; y evaluación final con las
métricas definidas en el marco teórico. El detalle de este procedimiento
se describe en el Capítulo IV, donde se presenta íntegramente en modo
proyectado.

\subsection{3.5. Criterio de Éxito: Umbral del 85\% y Plan de
Contingencia}\label{criterio-de-uxe9xito-umbral-del-85-y-plan-de-contingencia}

\subsubsection{3.5.1. Definición operacional del
umbral}\label{definiciuxf3n-operacional-del-umbral}

El \textbf{umbral del 85\%} se establece como \emph{criterio de éxito
metodológico} y se opera en dos vertientes:

\begin{itemize}
\tightlist
\item
  \textbf{Vertiente de red.} La celda experimental se considera que
  alcanza el criterio cuando la razón de entrega de tramas (PDR) es
  igual o superior al 85\%, con su intervalo de confianza
  correspondiente. La PDR se elige como métrica de referencia por ser la
  expresión directa de la confiabilidad entregada al usuario final y por
  ser una proporción directamente interpretable.
\item
  \textbf{Vertiente de modelado.} El modelo de Machine Learning se
  considera que alcanza el criterio cuando su desempeño en validación
  cruzada —exactitud y, de manera complementaria, el puntaje \emph{F}1—
  es igual o superior al 85\%.
\end{itemize}

La fijación del umbral \emph{antes} de la recolección es una decisión de
diseño metodológico: elimina la tentación de redefinir el criterio con
base en los resultados observados y convierte tanto el cumplimiento como
el incumplimiento en resultados interpretables. En consecuencia, la
investigación no se condiciona a que el 85\% se alcance: se condiciona a
que el 85\% se \emph{verifique con rigor} y a que, si no se alcanza, se
diagnostique técnicamente la causa mediante el procedimiento que sigue.

\subsubsection{3.5.2. Plan de contingencia de cinco
pasos}\label{plan-de-contingencia-de-cinco-pasos}

El plan se activa automáticamente cuando (i) ninguna celda del diseño
alcanza la PDR del 85\%, (ii) las celdas que lo alcanzan no son
reproducibles entre réplicas, o (iii) el modelo predictivo no alcanza la
exactitud del 85\% en validación cruzada. Sus cinco pasos son
secuenciales, excluyentes en su orden y cada uno posee un criterio de
cierre explícito.

\textbf{Paso 1 — Verificación del instrumental y de la trazabilidad.}
Acción: revisar la sincronización de relojes y su desfase registrado, la
integridad del firmware y su versión, la configuración vigente del punto
de acceso y del \emph{broker}, la continuidad del registro en la base de
datos y la correspondencia entre los tres registros (nodo, suscriptor y
\emph{broker}). Se repite una medición de referencia en un trayecto
conocido y se contrasta el RSSI con el instrumento independiente.
Criterio de cierre: se confirma que el instrumento produce registros
íntegros y consistentes, o se corrige la falla detectada y se reinicia
la campaña desde el punto afectado.

\textbf{Paso 2 — Re-aislamiento de causas exógenas.} Acción: verificar
mediante análisis espectral o revisión de la configuración la presencia
de redes vecinas, interferentes o actividad anómala no previstas;
repetir los ensayos afectados en horario de menor actividad y en canal
verificado; documentar las condiciones ambientales y de ocupación del
escenario. Criterio de cierre: se identifica y se acota o se documenta
como condición permanente del escenario la fuente exógena de
degradación.

\textbf{Paso 3 — Ampliación de la evidencia.} Acción: aumentar el número
de réplicas de las celdas infractoras, extender la duración de los
ensayos, revisar con criterios predefinidos —y no \emph{ad hoc}— la
posible exclusión de valores atípicos justificados documentalmente, y
evaluar la estabilidad de las estimaciones mediante el análisis de la
variabilidad entre réplicas. Criterio de cierre: se dispone de una
estimación con intervalo de confianza lo bastante estrecho como para
afirmar si el incumplimiento es real o producto de varianza muestral
insuficiente.

\textbf{Paso 4 — Reajuste de parámetros configurables del prototipo y
del protocolo.} Acción: modificar parámetros de configuración —no de
medición— con fundamento técnico: selección de canal Wi-Fi no solapado
distinto, ajuste de la potencia de transmisión si el hardware lo
permite, tamaño de carga útil, intervalo de \emph{keep alive},
temporizadores de sesión, profundidad de reintentos y nivel de servicio
MQTT. Se vuelve a ejecutar únicamente la submatriz de celdas afectadas.
Criterio de cierre: se identifica la configuración que alcanza el umbral
y se documenta como recomendación operativa, o se demuestra que ningún
ajuste de parámetros disponible permite alcanzarlo dentro de las
condiciones ensayadas.

\textbf{Paso 5 — Formalización de la desviación y re-umbralización
documentada.} Acción: si los pasos 1 a 4 no conducen al cumplimiento, se
emite un informe técnico de desviación que contiene: las celdas y
métricas afectadas, el umbral alcanzado con su intervalo de confianza,
el historial de acciones ejecutadas con su resultado, el diagnóstico
técnico de la causa y la propuesta de umbral alternativo debidamente
justificada. Los resultados se reportan íntegros, sin exclusión de
celdas ni modificación de datos recolectados; el umbral original se
conserva en la memoria del documento como referencia y el nuevo se
declara como \emph{objetivo condicional} para futuras iteraciones.
Criterio de cierre: el informe de desviación queda incorporado al
capítulo de resultados y a las recomendaciones, con decisión explícita
del asesor.

\textbf{Regla rectora del plan.} En ningún caso se modifica
retroactivamente el umbral sobre datos ya recolectados, no se eliminan
celdas del diseño sin registro documentado y no se declaran resultados
selectivamente: la integridad del protocolo experimental prevalece sobre
el cumplimiento del criterio.

{[}INSERTAR AQUÍ FLUJOGRAMA DE: Plan de contingencia de 5 pasos —
condiciones de activación, acciones y criterios de cierre de cada
paso{]}

\subsection{3.6. Alcance de la
Investigación}\label{alcance-de-la-investigaciuxf3n}

\textbf{Alcance por intervención.} La investigación \emph{diseña,
construye, instrumenta y valida} un prototipo de evaluación de QoS;
\emph{ejecuta} una campaña de mediciones estructurada bajo un diseño
factorial 3×3×3; \emph{analiza} los datos con métodos estadísticos y de
aprendizaje automático; y \emph{comunica} los resultados mediante
tableros. No se trata de una investigación orientada a la instalación
permanente de una red, sino a la validación de un método de medición y
de predicción.

\textbf{Alcance tecnológico.} Se evalúa exclusivamente el enlace Wi-Fi
IEEE 802.11 en la banda ISM de 2,4 GHz y el transporte MQTT sobre
TCP/IP. Quedan fuera de alcance las tecnologías de área personal (IEEE
802.15.4, Bluetooth Low Energy), las bandas de 5 y 6 GHz, los mecanismos
de planificación de capa física (MIMO de alto orden, modulation coding
scheme adaptativo gestionado por terceros), las redes mesh de salto
múltiple y los protocolos alternativos de transporte (CoAP, AMQP), cuyo
estudio se postula como línea futura.

\textbf{Alcance espacial.} Se limita a los interiores de la vivienda que
constituye el escenario experimental —una o dos plantas, con propagación
intra-planta—; no se evalúa penetración entre pisos ni escenarios
exteriores.

\textbf{Alcance de densidad.} La densidad de nodos se mantiene constante
y de valor moderado para aislar el efecto de los tres factores
ensayados; el efecto de la densidad como factor adicional queda
explícitamente fuera del alcance de esta iteración y se propone como
extensión del diseño a un factorial de cuatro factores.

\textbf{Alcance de seguridad y privacidad.} No se evalúa el desempeño de
los mecanismos de cifrado ni se incorporan ataques o condiciones
adversariales; se garantiza, en cambio, que la captura de datos se
limite a métricas de red y de radiofrecuencia, sin registro de
contenido.

\textbf{Alcance de viabilidad económica.} Se reporta la estimación de
costo por nodo y de la ausencia de infraestructura centralizada como
argumento de viabilidad, sin llevar a cabo un estudio de costo-vida útil
ni una comparación de mercado exhaustiva.

\textbf{Alcance temporal.} Corresponde al periodo de ejecución previsto
en el cronograma del plan de tesis, organizado en las cinco fases
descritas en el apartado 3.1.4.

\subsection{3.7. Propuesta de
Validación}\label{propuesta-de-validaciuxf3n}

La validación de la propuesta se estructura sobre cuatro mecanismos
complementarios que garantizan la rigurosidad científica de los
resultados:

\begin{enumerate}
\def\labelenumi{\arabic{enumi}.}
\tightlist
\item
  \textbf{Validación instrumental.} Contraste del RSSI reportado por el
  ESP32 contra un instrumento independiente; verificación de la
  sincronización de relojes; medición de referencia en trayecto conocido
  antes y después de la campaña. Toda discrepancia superior al
  tolerancia declarado se corrige antes de iniciar la recolección.
\item
  \textbf{Validación experimental.} Replicabilidad: cada celda se
  ejecuta en tres réplicas con orden aleatorizado, conforme al protocolo
  operativo del apartado 3.1.5; los tres registros (nodo, suscriptor,
  \emph{broker}) se triangulan para cada mensaje; las condiciones de
  ensayo se documentan de modo que una tercera parte pudiera reproducir
  la campaña.
\item
  \textbf{Validación estadística.} Verificación explícita de los
  supuestos del ANOVA, reporte de tamaños de efecto e intervalos de
  confianza, corrección por comparaciones múltiples y, para el modelado,
  partición entrenamiento–prueba con validación cruzada y fijación de
  semillas aleatorias para la reproducibilidad.
\item
  \textbf{Validación de utilidad.} Contrastación del instrumento frente
  a una línea base: el mismo enlace evaluado sin la instrumentación
  cruzada y las métricas que una operación convencional reportaría, de
  modo que se demuestre que el prototipo aporta información adicional
  utilizable. La capa de tableros se somete a verificación de
  coherencia: cada indicador visualizado debe poder rastrearse hasta la
  celda experimental y al registro original que lo sustenta.
\end{enumerate}

Adicionalmente, se prevé una \emph{revisión de hallazgos} con el asesor
antes del procesamiento final, con el fin de confirmar que la
interpretación de los resultados se mantiene dentro del marco hipotético
planteado y que el plan de contingencia, en caso de activarse, se
ejecuta conforme a los cinco pasos establecidos y en el orden fijado.

\section{CAPÍTULO IV: ANÁLISIS E INTERPRETACIÓN DE LOS
RESULTADOS}\label{capuxedtulo-iv-anuxe1lisis-e-interpretaciuxf3n-de-los-resultados}

\textbf{Modalidad de redacción.} El presente capítulo se formula
íntegramente en \emph{modo proyectado}: los verbos corresponden al
tiempo futuro y al condicional. La redacción en estos tiempos obedece a
una razón de rigor: la campaña de mediciones no se ha ejecutado y, por
tanto, no existen resultados que reportar. Este documento describe el
análisis \emph{previsto}, los procedimientos que \emph{se ejecutarán},
los criterios que \emph{se aplicarán} y las ramas de interpretación que
\emph{se activarían} según el comportamiento de los datos. En ningún
apartado se anticipan valores numéricos de medición, de desempeño de
modelos o de indicadores de tablero: toda cifra que aparezca
corresponderá exclusivamente a un umbral, a un nivel de factor o a un
parámetro metodológico fijado por adelantado en los capítulos
precedentes.

\subsection{4.1. Consideraciones Generales y Plan de
Análisis}\label{consideraciones-generales-y-plan-de-anuxe1lisis}

El plan de análisis seguirá la secuencia fijada en el protocolo
experimental: primero \emph{preparar} los datos, luego \emph{describir}
y \emph{verificar} su calidad estadística, después \emph{estimar los
efectos} del diseño factorial, a continuación \emph{construir y validar}
los modelos de aprendizaje automático, y finalmente \emph{comunicar} los
hallazgos mediante tableros y \emph{contrastarlos} con las preguntas de
investigación. Cada etapa posee un producto verificable y un criterio de
paso a la etapa siguiente.

El análisis se regirá por tres principios: (i) \emph{integridad} —no se
excluirán celdas ni observaciones sin registro documentado—; (ii)
\emph{reproducibilidad} —se fijarán semillas aleatorias, se versionará
el conjunto de datos y se conservará el \emph{pipeline} completo de
análisis—; y (iii) \emph{trazabilidad} —cada indicador publicado deberá
poder rastrearse hasta el registro original que lo sustenta. La base de
datos resultante de la campaña se denominará, de aquí en adelante, el
\emph{conjunto de mediciones}.

\subsection{4.2. Tratamiento y Preparación de los
Datos}\label{tratamiento-y-preparaciuxf3n-de-los-datos}

El conjunto de mediciones \emph{será} integrado desde tres registros
convergentes: el registro de publicación de los nodos emisores, el
registro de recepción de la estación suscriptora y el registro del
\emph{broker} MQTT. El primer paso \emph{consistirá} en la unión de los
tres conjuntos por el identificador único de mensaje y el número de
secuencia, operación que \emph{permitirá} discriminar la pérdida
ocurrida en el enlace inalámbrico de la pérdida ocurrida en la cola del
\emph{broker}.

Sobre la tabla unificada \emph{se ejecutarán} las siguientes
operaciones, en el orden indicado:

\begin{enumerate}
\def\labelenumi{\arabic{enumi}.}
\tightlist
\item
  \textbf{Verificación de integridad.} \emph{Se comprobará} la
  continuidad de los números de secuencia, la existencia de marcas de
  tiempo válidas, la ausencia de campos nulos en las variables críticas
  y la correspondencia entre el número de mensajes publicados y el
  número de registros almacenados. Los registros inconsistentes
  \emph{serán} marcados y, en su caso, \emph{serán} excluidos únicamente
  con causa documentada en el cuaderno de campo.
\item
  \textbf{Verificación de sincronización.} \emph{Se analizará} el
  desfase de los relojes informado por cada sesión de NTP; cuando el
  desfase \emph{supere} la tolerancia declarada en el protocolo, la
  sesión respectiva \emph{será} declarada inválida y sus ensayos
  \emph{se repetirán}.
\item
  \textbf{Construcción de variables derivadas.} \emph{Se calcularán} la
  latencia de un solo sentido y la latencia de ida y vuelta a partir de
  las marcas de tiempo, el jitter como desviación típica de las series
  móviles de retardo, la PDR y la pérdida como proporciones por celda y
  por réplica, el buen uso como carga útil por unidad de tiempo y los
  índices de radiofrecuencia como promedios y percentiles por ensayo.
\item
  \textbf{Codificación de factores.} Los escenarios, los niveles de
  carga y los niveles de QoS \emph{serán} codificados como variables
  categóricas ordenadas, y la etiqueta de celda factorial \emph{será}
  reconstruida de modo que cada observación \emph{quede} identificada
  con su combinación AiBjCk.
\item
  \textbf{Detección de anomalías.} \emph{Se aplicarán} criterios
  univariados predefinidos —basados en rangos intercuartílicos— para
  identificar valores atípicos; los valores detectados \emph{serán}
  inspeccionados contra el cuaderno de campo y \emph{se conservarán} en
  el conjunto de análisis salvo que exista evidencia documentada de
  falla instrumental.
\item
  \textbf{Estructuración para el modelado.} \emph{Se conformarán} la
  matriz de predictores —RSSI, SNR, nivel de ruido, distancia,
  escenario, carga y nivel de QoS— y las variables objetivo —estado de
  cumplimiento del umbral, latencia y jitter—, con una partición
  estratificada entre entrenamiento y prueba que \emph{se mantendrá}
  inmutable durante toda la fase de modelado.
\end{enumerate}

{[}INSERTAR AQUÍ FLUJOGRAMA DE: Flujo de tratamiento y preparación de
datos — unión de tres registros, verificación, variables derivadas,
codificación, detección de anomalías y partición{]}

\subsection{4.3. Análisis Descriptivo y Verificación de
Supuestos}\label{anuxe1lisis-descriptivo-y-verificaciuxf3n-de-supuestos}

\emph{Se calcularán} estadísticos descriptivos —media, mediana,
desviación típica, coeficiente de variación y percentiles 50, 95 y 99—
para cada métrica de QoS desagregada por factor y por celda. \emph{Se
construirán} histogramas y gráficos de caja por celda, que
\emph{permitirán} visualizar la asimetría de las distribuciones de
latencia, la presencia de colas largas asociadas a retransmisiones y la
dispersión de la PDR entre réplicas. \emph{Se elaborarán} mapas de calor
de PDR y de latencia media sobre las 27 celdas del diseño, cuya lectura
preliminar \emph{ofrecerá} una primera imagen del terreno de
cumplimiento del umbral.

Para cada variable respuesta \emph{se verificarán} los supuestos del
modelo de análisis de varianza: normalidad de los residuos mediante la
prueba de Shapiro-Wilk, homogeneidad de varianzas mediante la prueba de
Levene y aditividad del modelo. Cuando la normalidad \emph{resulte}
rechazada, \emph{se aplicarán} transformaciones —logarítmica para la
latencia, arcseno para las proporciones— y, si los supuestos
\emph{persisten} incumplidos, \emph{se recurrirá} a la alternativa no
paramétrica correspondiente, cuyo resultado \emph{se reportará} en
paralelo al del modelo paramétrico. La decisión tomada \emph{quedará}
documentada con su justificación antes de proceder a los contrastes.

{[}INSERTAR AQUÍ GRÁFICO: Comportamiento previsto de la latencia extremo
a extremo frente a la carga de ofrecimiento, diferenciado por nivel de
QoS MQTT{]}

{[}INSERTAR AQUÍ GRÁFICO: Mapa de calor previsto de PDR y latencia media
sobre las 27 celdas del diseño factorial{]}

\subsection{4.4. Análisis del Diseño Factorial
3×3×3}\label{anuxe1lisis-del-diseuxf1o-factorial-333}

\emph{Se ajustará} el modelo de análisis de varianza de tres vías con
interacciones para cada variable respuesta, conforme a la expresión
definida en el capítulo metodológico. La lectura de resultados
\emph{seguirá} un orden fijo:

\begin{enumerate}
\def\labelenumi{\arabic{enumi}.}
\tightlist
\item
  \emph{Se evaluará}, en primer lugar, la interacción de tres factores
  A×B×C; si \emph{resulta} significativa, la interpretación \emph{se
  hará} por capas de factores, ya que los efectos principales
  \emph{dejarían} de ser interpretables aisladamente.
\item
  \emph{Se estimarán} las interacciones de dos factores (A×B, A×C, B×C),
  con especial atención a A×B —la dependencia del efecto de la carga
  respecto del escenario de propagación— y a B×C —la dependencia del
  efecto de la carga respecto del nivel de servicio MQTT—, que
  corresponden al núcleo de la hipótesis de trabajo.
\item
  \emph{Se cuantificarán} los efectos principales de cada factor, con
  sus intervalos de confianza y sus tamaños de efecto.
\item
  \emph{Se aplicarán} contrastes de comparación múltiple por el método
  de Tukey para establecer, dentro de cada factor, qué niveles difieren
  significativamente.
\item
  \emph{Se calculará} la contribución porcentual de cada término a la
  varianza total, indicador que \emph{será} utilizado para ordenar los
  hallazgos por relevancia práctica y no únicamente por significancia
  estadística.
\end{enumerate}

Sobre la base del mismo modelo, \emph{se elaborarán} los intervalos de
confianza de la PDR por celda y \emph{se clasificarán} las 27 celdas en
dos categorías: aquellas cuyo intervalo \emph{se ubica} íntegramente por
encima del umbral del 85\%, aquellas cuyo intervalo \emph{se ubica}
íntegramente por debajo y aquellas cuyo intervalo \emph{lo atraviesa},
categoría que \emph{será} interpretada como evidencia insuficiente y, en
consecuencia, \emph{activará} el paso de ampliación de evidencia del
plan de contingencia.

{[}INSERTAR AQUÍ GRÁFICO: Diagrama de interacciones previsto — A×B, A×C,
B×C y contribución porcentual de cada término a la varianza total{]}

\subsection{4.5. Contrastación del Umbral del
85\%}\label{contrastaciuxf3n-del-umbral-del-85}

La contrastación \emph{se efectuará} en las dos vertientes definidas en
el capítulo metodológico. En la vertiente de red, \emph{se verificará}
si la PDR de cada celda alcanza o supera el 85\%; \emph{se determinará}
el número y la identidad de las celdas cumplidoras; y \emph{se
describirá} el patrón de distribución del cumplimiento a lo largo de los
factores. En la vertiente de modelado, \emph{se verificará} si el
desempeño del modelo en validación cruzada alcanza o supera el 85\% de
exactitud y, de manera complementaria, si el puntaje \emph{F}1 respalda
ese resultado.

Las ramas de decisión \emph{quedarán} formuladas de la siguiente manera:

\begin{itemize}
\tightlist
\item
  \textbf{Si} la vertiente de red alcanza el umbral en la mayoría de las
  celdas y el patrón \emph{resulta} coherente con lo previsto
  teóricamente —mejor desempeño en visibilidad directa, degradación
  creciente con la carga y diferencias ordenadas entre niveles de QoS—,
  \emph{se concluirá} que el prototipo permite verificar el criterio de
  éxito y \emph{se aceptará} la hipótesis de trabajo en esa dimensión.
\item
  \textbf{Si} el cumplimiento \emph{resulta} parcial y concentrado en
  condiciones específicas, \emph{se concluirá} que el umbral es
  alcanzable sólo bajo condiciones delimitadas, \emph{se identificarán}
  las fronteras operativas y \emph{se formularán} recomendaciones de
  configuración para cada escenario.
\item
  \textbf{Si} ninguna celda \emph{alcanza} el umbral, \emph{se activará}
  el plan de contingencia de cinco pasos en el orden establecido, y el
  análisis \emph{continuará} únicamente después de cerrar cada paso con
  su criterio correspondiente.
\item
  \textbf{Si} el modelo predictivo no alcanza la exactitud del 85\%,
  \emph{se discriminará} si la causa \emph{reside} en la insufuencia de
  predictores, en el desbalance de clases, en la varianza intrínseca de
  la respuesta o en la especificación del algoritmo, y \emph{se
  procederá} conforme a los pasos 3, 4 y 5 del plan de contingencia.
\end{itemize}

{[}INSERTAR AQUÍ FLUJOGRAMA DE: Árbol de decisión de contrastación del
umbral del 85\% — vertiente de red, vertiente de modelado y activación
del plan de contingencia{]}

\subsection{4.6. Modelado con
Scikit-learn}\label{modelado-con-scikit-learn}

La fase de modelado \emph{se ejecutará} en el entorno de Python
empleando la biblioteca Scikit-learn, siguiendo un flujo reproducible de
seis etapas.

\textbf{Etapa 1 — Partición y preprocesamiento.} El conjunto de datos
\emph{será} dividido en entrenamiento y prueba con estratificación
cuando la variable objetivo sea discreta; \emph{se aplicarán}
codificación de variables categóricas, imputación de ausentes y escalado
de variables continuas mediante canalizaciones (\emph{pipelines}) que
\emph{se ajustarán} exclusivamente sobre el conjunto de entrenamiento,
con el propósito de evitar la fuga de información.

\textbf{Etapa 2 — Selección de familias de modelos.} \emph{Se
entrenarán}, como mínimo, cuatro modelos: una regresión logística como
línea base interpretable, un árbol de decisión como modelo de referencia
no lineal, un bosque aleatorio como método de ensamble principal y una
regresión con regularización Ridge como línea base de la vertiente
continua. La selección \emph{no} se hará a priori: \emph{se compararán}
con las mismas métricas y el mismo esquema de validación.

\textbf{Etapa 3 — Ajuste de hiperparámetros.} \emph{Se buscarán} las
configuraciones óptimas mediante validación cruzada en \emph{k} pliegues
con búsqueda en malla o aleatorizada —número y profundidad de
estimadores, mínimo de muestras por hoja, parámetro de regularización—,
asegurando que la partición de prueba \emph{permanezca} intacta durante
toda la búsqueda.

\textbf{Etapa 4 — Evaluación.} Sobre el conjunto de prueba \emph{se
calcularán} exactitud, precisión, sensibilidad, especificidad, puntaje
\emph{F}1, área bajo la curva ROC y matriz de confusión para la
clasificación; y error absoluto medio, raíz del error cuadrático medio y
coeficiente de determinación para la regresión. Los resultados \emph{se
reportarán} con sus intervalos de confianza obtenidos por remuestreo.

\textbf{Etapa 5 — Interpretabilidad.} \emph{Se obtendrán} la importancia
de atributos del bosque aleatorio y los coeficientes de la regresión
regularizada; \emph{se contrastarán} ambos criterios; y \emph{se
verificará} si los predictores de radiofrecuencia y de configuración
—RSSI, SNR, distancia, carga y nivel de QoS— figuran entre los más
relevantes. Este paso \emph{será} el que vincule el modelo con el marco
físico de la investigación: un modelo con desempeño elevado pero cuyos
predictores dominantes \emph{sean} ajenos al problema planteado
\emph{será} considerado insuficiente para los fines de la tesis.

\textbf{Etapa 6 — Comparación y selección.} \emph{Se elegirá} el modelo
que ofrezca el mejor compromiso entre desempeño, estabilidad entre
pliegues e interpretabilidad, y \emph{se documentarán} las razones de la
selección. \emph{Se conservarán} el modelo serializado, la semilla
aleatoria y la configuración final, de modo que su evaluación
\emph{pueda} repetirse de forma idéntica.

\textbf{Ramas condicionales.} Si el modelo seleccionado alcanza la
exactitud y el puntaje \emph{F}1 del 85\% en validación cruzada y en el
conjunto de prueba, \emph{se concluirá} que la correlación entre
variables físicas de radiofrecuencia y variables dinámicas de red
\emph{es} modelable con precisión suficiente para fines de evaluación de
QoS, y \emph{se aceptará} la hipótesis en esa dimensión. Si el desempeño
\emph{resulta} inferior, \emph{se analizarán} las causas técnicas
—predictores insuficientes, ruido de medición, efecto de factores no
controlados, tamaño muestral, desbalance de clases—, \emph{se aplicarán}
los ajustes correspondientes y, en última instancia, \emph{se reportará}
el desempeño alcanzado sin forzar el cumplimiento. Si la importancia de
predictores \emph{mostrara} que el nivel de QoS o la carga dominan al
RSSI, \emph{se concluirá} que el comportamiento de la red interior en
las condiciones ensayadas está gobernado más por la capa de aplicación y
de acceso que por la atenuación, hallazgo que \emph{sería} coherente con
la hipótesis de conflicto estructural entre capas planteada en la
introducción.

{[}INSERTAR AQUÍ FLUJOGRAMA DE: Pipeline de Machine Learning con
Scikit-learn — 6 etapas: partición, selección de familias, ajuste de
hiperparámetros, evaluación, interpretabilidad y selección final{]}

\subsection{4.7. Visualización y Tableros en Power
BI}\label{visualizaciuxf3n-y-tableros-en-power-bi}

Los resultados \emph{serán} trasladados a Power BI mediante la
importación del conjunto de mediciones y de las tablas resumen
producidas por el análisis estadístico y por las predicciones del
modelo.

\textbf{Modelado de datos.} \emph{Se construirá} un modelo estrella: una
tabla de hechos con una fila por mensaje o por ensayo, y tablas de
dimensión por escenario de propagación, nivel de carga, nivel de QoS,
celda factorial y calendario. \emph{Se definirán} relaciones unívocas
entre hechos y dimensiones y \emph{se verificará} la cardinalidad antes
de publicar.

\textbf{Medidas y lógica de negocio.} \emph{Se crearán} medidas de DAX
para la PDR, la pérdida, la latencia media y percentil 95, el jitter, el
buen uso y el porcentaje de celdas que cumplen el umbral, incluyendo el
cálculo condicional de cumplimiento respecto del 85\% y la variación de
cada indicador respecto de la línea base. Todas las medidas \emph{serán}
validadas mediante contraste manual contra los valores calculados
directamente en el entorno de análisis.

\textbf{Tableros previstos.} \emph{Se diseñarán} al menos tres tableros:
(i) un \emph{tablero de estado}, con semáforos de cumplimiento del
umbral por celda y por métrica; (ii) un \emph{tablero de diagnóstico},
con series temporales de latencia y jitter, mapas de calor por escenario
y distribuciones de RSSI; y (iii) un \emph{tablero predictivo}, con la
comparación entre los valores medidos y los valores estimados por el
modelo, la importancia de predictores y el porcentaje de aciertos de
clasificación. \emph{Se configurarán} filtros cruzados por escenario,
carga y nivel de QoS, así como la posibilidad de profundizar
(\emph{drill-through}) desde el indicador agregado hasta el registro
individual.

\textbf{Verificación de coherencia.} Antes de la publicación, \emph{se
contrastará} cada indicador visualizado contra el valor calculado en el
procesamiento estadístico; cualquier discrepancia \emph{impedirá} la
publicación del tablero hasta ser resuelta. Los tableros \emph{quedarán}
enlazados al conjunto de mediciones versionado, de modo que su
actualización \emph{sea} reproducible ante la adición de nuevas
campañas.

{[}INSERTAR AQUÍ TABLA: Modelo estrella y esquema de relaciones para los
tableros de control en Power BI{]}

{[}INSERTAR AQUÍ IMAGEN TÉCNICA: Esquema de los tableros de Power BI —
estado, diagnóstico y predictivo, con filtros cruzados y
drill-through{]}

\subsection{4.8. Interpretación y Contrastación con las Preguntas de
Investigación}\label{interpretaciuxf3n-y-contrastaciuxf3n-con-las-preguntas-de-investigaciuxf3n}

La interpretación \emph{se organizará} siguiendo el sistema de preguntas
formulado en el capítulo introductorio:

\begin{itemize}
\tightlist
\item
  \textbf{Respecto de la PE1.} \emph{Se evaluará} el ajuste del modelo
  de pérdida de trayecto calibrado al escenario mediante el error medio
  cuadrático y el coeficiente de determinación, y \emph{se discutirá} la
  correspondencia entre el RSSI observado y la condición de obstrucción
  de cada escenario. Si el modelo \emph{resulta} adecuado, \emph{se
  afirmará} que la familia de variables físicas del instrumento es
  representativa del canal; de lo contrario, \emph{se explicitarán} las
  fuentes de dispersión no capturadas —multipropagación, ocupación del
  local, interferencia vecina— y su efecto sobre el análisis posterior.
\item
  \textbf{Respecto de la PE2.} \emph{Se interpretarán} los efectos
  principales e interacciones del factorial a la luz de los marcos
  teóricos del acceso al medio y de la semántica de confirmación de
  MQTT: \emph{se examinará}, por ejemplo, si la sobrecarga asociada al
  nivel QoS 2 \emph{se incrementa} con la carga y si esa dependencia
  \emph{se atenúa} o \emph{se agrava} según el escenario de propagación.
  Los hallazgos \emph{se contrastarán} con el ordenamiento cualitativo
  reportado en el estado del arte, \emph{se señalarán} las coincidencias
  y las divergencias y \emph{se argumentarán} estas últimas a partir de
  las diferencias de entorno y de carga.
\item
  \textbf{Respecto de la PE3.} \emph{Se discutirá} el desempeño relativo
  de los algoritmos evaluados, la estabilidad entre pliegues, la
  estructura de importancia de predictores y la capacidad del modelo
  para discriminar el cumplimiento del umbral. \emph{Se relacionará} la
  magnitud de la contribución de los predictores de radiofrecuencia con
  la importancia que el estado del arte les atribuye en contextos
  distintos.
\item
  \textbf{Respecto de la PE4.} \emph{Se evaluará} la utilidad del
  instrumento a partir de la cobertura de información que aporta
  respecto de la operación convencional, la trazabilidad de los tableros
  y la estimación de costo por nodo, y \emph{se situará} el resultado en
  una escala explícita de madurez tecnológica.
\end{itemize}

La contrastación final \emph{incluirá} una tabla de correspondencia
entre cada pregunta, cada objetivo específico, la evidencia producida y
la conclusión alcanzada, de modo que la coherencia interna del trabajo
\emph{pueda} verificarse de un solo vistazo.

{[}INSERTAR AQUÍ TABLA: Matriz de correspondencia PE1–PE4 — pregunta,
objetivo específico, evidencia producida y conclusión prevista{]}

\subsection{4.9. Validación, Límites y Amenazas a la
Validez}\label{validaciuxf3n-luxedmites-y-amenazas-a-la-validez}

\emph{Se someterá} el conjunto de resultados a cuatro controles de
validez:

\begin{itemize}
\tightlist
\item
  \textbf{Validez interna.} \emph{Se verificará} que las diferencias
  observadas \emph{sean} atribuibles a los factores ensayados,
  descartando por registro documentado interferencias exógenas no
  controladas; \emph{se revisará} si el orden aleatorio de los ensayos
  \emph{neutralizó} posibles tendencias temporales.
\item
  \textbf{Validez externa.} \emph{Se discutirá} cuidadosamente la
  transferibilidad de los hallazgos a escenarios distintos del ensayado,
  considerando la banda de operación, la densidad fija de nodos, la
  geometría de las instalaciones y el protocolo de transporte evaluado.
\item
  \textbf{Validez de constructo.} \emph{Se argumentará} en qué medida la
  PDR y la exactitud del modelo \emph{representan} la noción de calidad
  de servicio definida en el marco teórico, y \emph{se reconocerán} los
  aspectos de la experiencia de calidad que el instrumento no captura
  —carga de trabajo del usuario, semántica de la aplicación, seguridad—.
\item
  \textbf{Fiabilidad.} \emph{Se reportarán} la repetibilidad de las
  métricas entre réplicas, la consistencia entre los tres registros
  convergentes y la reproducibilidad del análisis completo a partir del
  conjunto de datos versionado.
\end{itemize}

\textbf{Límites previstos.} \emph{Se anticipa} que los resultados quedan
delimitados a la banda de 2,4 GHz, a una densidad de nodos constante y a
interiores de una planta; que la interferencia de redes vecinas no
controladas \emph{podría} introducir varianza no atribuible a los
factores; y que el tamaño muestral \emph{podría} limitar la precisión de
las estimaciones en celdas de respuesta dispersa. Estas amenazas
\emph{serán} tratadas explícitamente en las conclusiones, con indicación
de su dirección probable —sesgo conservador o sesgo optimista— sobre los
hallazgos.

\textbf{Proyección.} En función de los resultados obtenidos, \emph{se
formularán} recomendaciones de configuración operativa para redes
interiores —escenario, carga y nivel de QoS recomendados por condición—,
\emph{se propondrán} extensiones del diseño —incorporación de la
densidad de nodos como cuarto factor, evaluación en banda de 5 GHz,
incorporación de protocolos alternativos de transporte— y \emph{se
definirá} la ruta para elevar el prototipo a un grado de madurez
tecnológica superior mediante pruebas en instalaciones y condiciones de
operación reales.

\end{document}
